% Môn: Vật lý
% Lớp: 11
% Chương: Chương II. Sóng
% Bài: L11C2 Bài 11. Sóng điện từ
% Số câu: 168
% App đọc file này trực tiếp — sửa rồi Commit trên GitHub, bấm Đồng bộ GitHub .tex

% L11C2 Bài 11. Sóng điện từ

% ===== Câu 1 =====
% ID: L11C2B11-01-DS
% Mức: TH
\dangbt{Bài toán truyền sóng điện từ qua các môi trường}
\begin{ex}
% ID: L11C2B11-01-DS
% Mức: TH
Trong chuyên đề “Bài toán truyền sóng điện từ qua các môi trường”, Sóng điện từ có $f=150$ MHz trong chân không.
\choiceTF
{\True Bước sóng bằng $2\,\text{m}$.}
{\True Chu kì bằng $6,67$ ns.}
{\True Trong 1 micro giây sóng truyền $300\,\text{m}$.}
{Đổi sang môi trường khác làm tần số giảm.}
\loigiai{
\begin{itemchoice}
\item (Đúng) Bước sóng bằng $2\,\text{m}$. (Vì): Bước sóng bằng $2\,\text{m}$. b) Đ: Chu kì bằng $6,67$ ns. c) Đ: Trong 1 micro giây sóng truyền $300\,\text{m}$. d) S: Đổi sang môi trường khác làm tần số giảm
\item (Đúng) Chu kì bằng $6,67$ ns.
\item (Đúng) Trong 1 micro giây sóng truyền $300\,\text{m}$.
\item (Sai) Đổi sang môi trường khác làm tần số giảm.
\end{itemchoice}
}
\end{ex}

% ===== Câu 2 =====
% ID: L11C2B11-01-TL
% Mức: TH
\begin{ex}
% ID: L11C2B11-01-TL
% Mức: TH
Trong chuyên đề “Bài toán truyền sóng điện từ qua các môi trường”, Một đài phát ở tần số 100 MHz. Tính bước sóng trong chân không và trong môi trường có tốc độ $2,4\times10^8$ m/s.
\choiceTF

\loigiai{
Trong chân không $\lambda=3$ m; trong môi trường $\lambda=2,4$ m; tần số không đổi.
}
\end{ex}

% ===== Câu 3 =====
% ID: L11C2B11-01-TLN
% Mức: TH
\begin{ex}
% ID: L11C2B11-01-TLN
% Mức: TH
Trong chuyên đề “Bài toán truyền sóng điện từ qua các môi trường”, Sóng điện từ 60 MHz truyền trong chân không. Bước sóng tính theo m bằng bao nhiêu?
\shortans{5}
\loigiai{
Xác định đúng dữ kiện và đơn vị, áp dụng hệ thức của dạng bài rồi tính được kết quả 5.
}
\end{ex}

% ===== Câu 4 =====
% ID: L11C2B11-01-TN
% Mức: NB
\begin{ex}
% ID: L11C2B11-01-TN
% Mức: NB
Trong chuyên đề “Bài toán truyền sóng điện từ qua các môi trường”, Sóng điện từ trong chân không có tần số 100 MHz. Bước sóng bằng
\choice
{$0,3\,\text{m}$}
{$30\,\text{m}$}
{$300\,\text{m}$}
{\True $3\,\text{m}$}
\loigiai{
Đáp án D. $3\,\text{m}$. Đối chiếu định nghĩa, hệ thức hoặc dữ kiện của bài toán cho kết luận trên.
}
\end{ex}

% ===== Câu 5 =====
% ID: L11C2B11-02-TL
% Mức: TH
\begin{ex}
% ID: L11C2B11-02-TL
% Mức: TH
Trong chuyên đề “Bài toán truyền sóng điện từ qua các môi trường”, Ánh sáng đỏ có bước sóng 750 nm, ánh sáng tím 400 nm. Tính tần số mỗi ánh sáng và so sánh.
\choiceTF

\loigiai{
$f_đ=4,0\times10^{14}$ Hz; $f_t=7,5\times10^{14}$ Hz; ánh sáng tím có tần số lớn hơn.
}
\end{ex}

% ===== Câu 6 =====
% ID: L11C2B11-02-TLN
% Mức: TH
\begin{ex}
% ID: L11C2B11-02-TLN
% Mức: TH
Trong chuyên đề “Bài toán truyền sóng điện từ qua các môi trường”, Bức xạ có bước sóng 400 nm trong chân không. Tần số tính theo $10^{14}$ Hz bằng bao nhiêu?
\shortans{07/05/2026}
\loigiai{
Xác định đúng dữ kiện và đơn vị, áp dụng hệ thức của dạng bài rồi tính được kết quả 7.5.
}
\end{ex}

% ===== Câu 7 =====
% ID: L11C2B11-03-DS
% Mức: TH
\begin{ex}
% ID: L11C2B11-03-DS
% Mức: TH
Trong chuyên đề “Bài toán truyền sóng điện từ qua các môi trường”, Sóng 75 MHz truyền trong môi trường với tốc độ $2,25\times10^8$ m/s.
\choiceTF
{\True Bước sóng bằng $3\,\text{m}$.}
{\True Chu kì bằng $1,33\times10^{-8}$ s.}
{\True Tần số không đổi khi truyền sang chân không.}
{Bước sóng trong chân không vẫn bằng $3\,\text{m}$.}
\loigiai{
\begin{itemchoice}
\item (Đúng) Bước sóng bằng $3\,\text{m}$. (Vì): Bước sóng bằng $3\,\text{m}$. b) Đ: Chu kì bằng $1,33\times10^{-8}$ s. c) Đ: Tần số không đổi khi truyền sang chân không. d) S: Bước sóng trong chân không vẫn bằng $3\,\text{m}$
\item (Đúng) Chu kì bằng $1,33\times10^{-8}$ s.
\item (Đúng) Tần số không đổi khi truyền sang chân không.
\item (Sai) Bước sóng trong chân không vẫn bằng $3\,\text{m}$.
\end{itemchoice}
}
\end{ex}

% ===== Câu 8 =====
% ID: L11C2B11-03-TL
% Mức: VD
\begin{ex}
% ID: L11C2B11-03-TL
% Mức: VD
Trong chuyên đề “Bài toán truyền sóng điện từ qua các môi trường”, Một xung radar đi tới mục tiêu rồi phản xạ về sau 200 micro giây. Tính khoảng cách tới mục tiêu.
\choiceTF

\loigiai{
Quãng đường hai chiều $ct$ nên $d=ct/2=3\times10^8\times200\times10^{-6}/2=30000$ m = $30\,\text{km}$.
}
\end{ex}

% ===== Câu 9 =====
% ID: L11C2B11-03-TLN
% Mức: TH
\begin{ex}
% ID: L11C2B11-03-TLN
% Mức: TH
Trong chuyên đề “Bài toán truyền sóng điện từ qua các môi trường”, Sóng có tần số 80 MHz truyền với tốc độ $2,4\times10^8$ m/s. Bước sóng tính theo m bằng bao nhiêu?
\shortans{3}
\loigiai{
Xác định đúng dữ kiện và đơn vị, áp dụng hệ thức của dạng bài rồi tính được kết quả 3.
}
\end{ex}

% ===== Câu 10 =====
% ID: L11C2B11-04-DS
% Mức: VD
\begin{ex}
% ID: L11C2B11-04-DS
% Mức: VD
Trong chuyên đề “Bài toán truyền sóng điện từ qua các môi trường”, Xét các nhận định tổng hợp thuộc dạng “Bài toán truyền sóng điện từ qua các môi trường”.
\choiceTF
{\True $3\,\text{m}$.}
{\True $5\times10^{14}$ Hz.}
{\True $4\,\text{m}$.}
{$0,3\,\text{m}$.}
\loigiai{
\begin{itemchoice}
\item (Đúng) $3\,\text{m}$. (Vì): $3\,\text{m}$. b) Đ: $5\times10^{14}$ Hz. c) Đ: $4\,\text{m}$. d) S: $0,3\,\text{m}$
\item (Đúng) $5\times10^{14}$ Hz.
\item (Đúng) $4\,\text{m}$.
\item (Sai) $0,3\,\text{m}$.
\end{itemchoice}
}
\end{ex}

% ===== Câu 11 =====
% ID: L11C2B11-04-TL
% Mức: VD
\begin{ex}
% ID: L11C2B11-04-TL
% Mức: VD
Trong chuyên đề “Bài toán truyền sóng điện từ qua các môi trường”, Phân tích các bước lập luận và chỉ ra điều kiện áp dụng trong bài toán: Một đài phát ở tần số 100 MHz. Tính bước sóng trong chân không và trong môi trường có tốc độ $2,4\times10^8$ m/s.
\choiceTF

\loigiai{
Trong chân không $\lambda=3$ m; trong môi trường $\lambda=2,4$ m; tần số không đổi. Cần nêu rõ điều kiện áp dụng, đổi đơn vị thống nhất và kiểm tra ý nghĩa vật lí của kết quả.
}
\end{ex}

% ===== Câu 12 =====
% ID: L11C2B11-04-TLN
% Mức: VD
\begin{ex}
% ID: L11C2B11-04-TLN
% Mức: VD
Trong chuyên đề “Bài toán truyền sóng điện từ qua các môi trường”, Từ kết quả của tình huống sau, nếu đại lượng cần tìm được nhân với hệ số 2 thì giá trị mới bằng bao nhiêu? Sóng điện từ 60 MHz truyền trong chân không. Bước sóng tính theo m bằng bao nhiêu?
\shortans{10}
\loigiai{
Xác định đúng dữ kiện và đơn vị, áp dụng hệ thức của dạng bài rồi tính được kết quả 10.
}
\end{ex}

% ===== Câu 13 =====
% ID: L11C2B11-04-TN
% Mức: TH
\begin{ex}
% ID: L11C2B11-04-TN
% Mức: TH
Trong chuyên đề “Bài toán truyền sóng điện từ qua các môi trường”, Sóng điện từ 60 MHz truyền trong chân không. Bước sóng tính theo m bằng bao nhiêu? Chọn giá trị đúng.
\choice
{02/05/2026}
{6}
{\True 5}
{10}
\loigiai{
Đáp án C. 5. Đối chiếu định nghĩa, hệ thức hoặc dữ kiện của bài toán cho kết luận trên.
}
\end{ex}

% ===== Câu 14 =====
% ID: L11C2B11-05-DS
% Mức: VD
\begin{ex}
% ID: L11C2B11-05-DS
% Mức: VD
Trong chuyên đề “Bài toán truyền sóng điện từ qua các môi trường”, Đánh giá cách xử lí các dữ kiện định lượng của dạng “Bài toán truyền sóng điện từ qua các môi trường”.
\choiceTF
{\True Kết quả của bài toán thứ nhất là 5.}
{\True Kết quả của bài toán thứ hai là 7.5.}
{\True Kết quả của bài toán thứ ba là 3.}
{Có thể bỏ qua hoàn toàn đơn vị và điều kiện vật lí khi tính toán.}
\loigiai{
\begin{itemchoice}
\item (Đúng) Kết quả của bài toán thứ nhất là 5. (Vì): Kết quả của bài toán thứ nhất là 5. b) Đ: Kết quả của bài toán thứ hai là 7.5. c) Đ: Kết quả của bài toán thứ ba là 3. d) S: Có thể bỏ qua hoàn toàn đơn vị và điều kiện vật lí khi tính toán
\item (Đúng) Kết quả của bài toán thứ hai là 7.5.
\item (Đúng) Kết quả của bài toán thứ ba là 3.
\item (Sai) Có thể bỏ qua hoàn toàn đơn vị và điều kiện vật lí khi tính toán.
\end{itemchoice}
}
\end{ex}

% ===== Câu 15 =====
% ID: L11C2B11-05-TL
% Mức: VD
\begin{ex}
% ID: L11C2B11-05-TL
% Mức: VD
Trong chuyên đề “Bài toán truyền sóng điện từ qua các môi trường”, Một học sinh giải bài sau nhưng bỏ qua điều kiện vật lí. Hãy sửa cách làm và trình bày lời giải đúng: Ánh sáng đỏ có bước sóng 750 nm, ánh sáng tím 400 nm. Tính tần số mỗi ánh sáng và so sánh.
\choiceTF

\loigiai{
Cách giải đúng: $f_đ=4,0\times10^{14}$ Hz; $f_t=7,5\times10^{14}$ Hz; ánh sáng tím có tần số lớn hơn. Sai lầm cần tránh là dùng hệ thức khi chưa xác định điều kiện biên, môi trường hoặc đơn vị.
}
\end{ex}

% ===== Câu 16 =====
% ID: L11C2B11-05-TLN
% Mức: VD
\begin{ex}
% ID: L11C2B11-05-TLN
% Mức: VD
Trong chuyên đề “Bài toán truyền sóng điện từ qua các môi trường”, Từ kết quả của tình huống sau, nếu đại lượng cần tìm được nhân với hệ số 0.5 thì giá trị mới bằng bao nhiêu? Bức xạ có bước sóng 400 nm trong chân không. Tần số tính theo $10^{14}$ Hz bằng bao nhiêu?
\shortans{3.75}
\loigiai{
Xác định đúng dữ kiện và đơn vị, áp dụng hệ thức của dạng bài rồi tính được kết quả 3.75.
}
\end{ex}

% ===== Câu 17 =====
% ID: L11C2B11-05-TN
% Mức: TH
\begin{ex}
% ID: L11C2B11-05-TN
% Mức: TH
Trong chuyên đề “Bài toán truyền sóng điện từ qua các môi trường”, Bức xạ có bước sóng 400 nm trong chân không. Tần số tính theo $10^{14}$ Hz bằng bao nhiêu? Chọn giá trị đúng.
\choice
{15}
{3.75}
{08/05/2026}
{\True 07/05/2026}
\loigiai{
Đáp án D. 7.5. Đối chiếu định nghĩa, hệ thức hoặc dữ kiện của bài toán cho kết luận trên.
}
\end{ex}

% ===== Câu 18 =====
% ID: L11C2B11-06-TL
% Mức: VD
\begin{ex}
% ID: L11C2B11-06-TL
% Mức: VD
Trong chuyên đề “Bài toán truyền sóng điện từ qua các môi trường”, Mở rộng tình huống sau thành một ví dụ thực tế và trình bày phương pháp giải: Một xung radar đi tới mục tiêu rồi phản xạ về sau 200 micro giây. Tính khoảng cách tới mục tiêu.
\choiceTF

\loigiai{
Quãng đường hai chiều $ct$ nên $d=ct/2=3\times10^8\times200\times10^{-6}/2=30000$ m = $30\,\text{km}$. Có thể kiểm chứng bằng cách thay kết quả vào dữ kiện ban đầu và đối chiếu với hiện tượng thực tế.
}
\end{ex}

% ===== Câu 19 =====
% ID: L11C2B11-06-TLN
% Mức: VD
\begin{ex}
% ID: L11C2B11-06-TLN
% Mức: VD
Trong chuyên đề “Bài toán truyền sóng điện từ qua các môi trường”, Từ kết quả của tình huống sau, nếu đại lượng cần tìm được nhân với hệ số 1.5 thì giá trị mới bằng bao nhiêu? Sóng có tần số 80 MHz truyền với tốc độ $2,4\times10^8$ m/s. Bước sóng tính theo m bằng bao nhiêu?
\shortans{04/05/2026}
\loigiai{
Xác định đúng dữ kiện và đơn vị, áp dụng hệ thức của dạng bài rồi tính được kết quả 4.5.
}
\end{ex}

% ===== Câu 20 =====
% ID: L11C2B11-06-TN
% Mức: TH
\begin{ex}Một sóng điện từ có tần số $80$ MHz đang lan truyền trong một môi trường với tốc độ $2,4 \times 10^8$ m/s. Hãy xác định bước sóng của sóng này.
\choice
{\True 3 m}
{6 m}
{2 m}
{4 m}
\loigiai{
Để tính bước sóng $\lambda$, ta sử dụng công thức $\lambda = \frac{v}{f}$, trong đó $v$ là tốc độ truyền sóng và $f$ là tần số của sóng.
Theo đề bài, ta có:
$f = 80$ MHz $= 80 \times 10^6$ Hz
$v = 2,4 \times 10^8$ m/s
Thay các giá trị vào công thức, ta được:
$\lambda = \frac{2,4 \times 10^8}{80 \times 10^6} = \frac{2,4 \times 10^2}{80} = \frac{240}{80} = 3$ m.
Vậy, bước sóng của sóng điện từ là $3$ m. Chọn phương án A.
}
\end{ex}

% ===== Câu 21 =====
% ID: L11C2B11-07-DS
% Mức: TH
\dangbt{Nhận biết bản chất và tính chất của sóng điện từ}
\begin{ex}
% ID: L11C2B11-07-DS
% Mức: TH
Sóng điện từ đi từ chân không vào điện môi.
\choiceTF
{\True Tần số không đổi.}
{\True Tốc độ thường giảm.}
{\True Bước sóng giảm nếu tốc độ giảm.}
{Chu kì tăng gấp đôi trong mọi điện môi.}
\loigiai{
\begin{itemchoice}
\item (Đúng) Tần số không đổi. (Vì): Tần số không đổi. b) Đ: Tốc độ thường giảm. c) Đ: Bước sóng giảm nếu tốc độ giảm. d) S: Chu kì tăng gấp đôi trong mọi điện môi
\item (Đúng) Tốc độ thường giảm.
\item (Đúng) Bước sóng giảm nếu tốc độ giảm.
\item (Sai) Chu kì tăng gấp đôi trong mọi điện môi.
\end{itemchoice}
}
\end{ex}

% ===== Câu 22 =====
% ID: L11C2B11-07-TL
% Mức: TH
\dangbt{Chưa có dạng}
\begin{ex}
% ID: L11C2B11-07-TL
% Mức: TH
Sóng vô tuyến ngắn có thể được sử dụng để đo khoảng cách từ Trái Đất đến Mặt Trăng, bằng cách phát một tín hiệu từ Trái Đất tới Mặt Trăng và thu tín hiệu trở lại, đo khoảng thời gian từ khi phát đến khi nhận tín hiệu. Khoảng thời gian từ khi phát tới khi nhận được tín hiệu trở lại là $2,5\,\mathrm{s}$. Biết tốc độ của sóng vô tuyến này là $3.108 \mathrm{m/s}$ và có tần số $107 \mathrm{Hz}$. Tính: a) Khoảng cách từ Mặt Trăng tới Trái Đất. b) Bước sóng của sóng vô tuyến đã sử dụng. $\nguon{ly11 kntt.pdf}$
\choiceTF

\loigiai{
a) Khoảng cách từ Trái Đất tới Mặt Trăng: $\text{d} = \dfrac{\text{ct}}{2} = \dfrac{3 \cdot 10^{8} \cdot 2,5}{2} = 3,75 \cdot 10^{8}\text{m}$ b) $\lambda = \dfrac{c}{f} = \dfrac{3 \cdot 10^{8}}{10^{7}} = 30\text{m}$.
}
\end{ex}

% ===== Câu 23 =====
% ID: L11C2B11-07-TLN
% Mức: TH
\dangbt{Nhận biết bản chất và tính chất của sóng điện từ}
\begin{ex}
% ID: L11C2B11-07-TLN
% Mức: TH
Sóng điện từ truyền trong chân không $0,002\,\text{s}$. Quãng đường tính theo km bằng bao nhiêu?
\shortans{600}
\loigiai{
Xác định đúng dữ kiện và đơn vị, áp dụng hệ thức của dạng bài rồi tính được kết quả 600.
}
\end{ex}

% ===== Câu 24 =====
% ID: L11C2B11-07-TN
% Mức: VD
\dangbt{Bài toán truyền sóng điện từ qua các môi trường}
\begin{ex}Một đài phát sóng điện từ hoạt động ở tần số $100$ MHz. Hãy xác định bước sóng của sóng này trong chân không và trong một môi trường vật chất có tốc độ truyền sóng là $2,4 \times 10^8$ m/s.
\choice
{Bước sóng của sóng điện từ không phụ thuộc vào môi trường truyền sóng.}
{\True Trong chân không, bước sóng của sóng điện từ là $2,4§ m}
{Để giải bài toán này, không cần sử dụng các định nghĩa hay dữ kiện đã cho.}
{Bước sóng của sóng điện từ luôn bằng $0$ trong mọi trường hợp.}
\loigiai{
Tần số của sóng là $f = 100$ MHz $= 100 \times 10^6$ Hz $= 10^8$ Hz.
Trong chân không, tốc độ truyền sóng điện từ là $c = 3 \times 10^8$ m/s.
Bước sóng trong chân không là $\lambda_0 = \frac{c}{f} = \frac{3 \times 10^8}{10^8} = 3$ m.
Trong môi trường vật chất có tốc độ truyền sóng $v = 2,4 \times 10^8$ m/s.
Bước sóng trong môi trường này là $\lambda = \frac{v}{f} = \frac{2,4 \times 10^8}{10^8} = 2,4$ m.
Đối chiếu với các phương án, phương án B là đúng.
}
\end{ex}

% ===== Câu 25 =====
% ID: L11C2B11-08-DS
% Mức: TH
\dangbt{Nhận biết bản chất và tính chất của sóng điện từ}
\begin{ex}
% ID: L11C2B11-08-DS
% Mức: TH
Các hiện tượng của sóng điện từ.
\choiceTF
{\True Có thể phản xạ.}
{\True Có thể khúc xạ.}
{\True Có thể giao thoa và nhiễu xạ.}
{Không bao giờ truyền qua chân không.}
\loigiai{
\begin{itemchoice}
\item (Đúng) Có thể phản xạ. (Vì): Có thể phản xạ. b) Đ: Có thể khúc xạ. c) Đ: Có thể giao thoa và nhiễu xạ. d) S: Không bao giờ truyền qua chân không
\item (Đúng) Có thể khúc xạ.
\item (Đúng) Có thể giao thoa và nhiễu xạ.
\item (Sai) Không bao giờ truyền qua chân không.
\end{itemchoice}
}
\end{ex}

% ===== Câu 26 =====
% ID: L11C2B11-08-TL
% Mức: TH
\dangbt{Chưa có dạng}
\begin{ex}
% ID: L11C2B11-08-TL
% Mức: TH
*. Một anten radar phát ra những sóng điện từ đến vật đang chuyển động về phía radar. Thời gian từ lúc anten phát sóng đến lúc nhận sóng phản xạ từ vật trở lại là 80 $\mus$. Sau hai phút, đo lần thứ hai, thời gian từ lúc phát đến lúc nhận lần này là 76 $\mus$. Tính tốc độ trung bình của vật. Coi tốc độ c ủa sóng điện từ trong không khí bằng 3 $\cdot108 \mathrm{m/s}$. $\nguon{ly11 kntt.pdf}$
\choiceTF

\loigiai{
Lần 1: $d_{1} = \dfrac{ct_{1}}{2} = \dfrac{3 \cdot 10^{8} \cdot 80 \cdot 10^{-6}}{2}=12000\,\mathrm{m}$.

Lần 2: $\text{d}_{2} = \dfrac{\text{ct}_{2}}{2} = \dfrac{3 \cdot 10^{8} \cdot 76 \cdot 10^{-6}}{2}$ = $11400\,\mathrm{m}$.

$\left. \Rightarrow\overline{v} = \dfrac{d_{2} - d_{1}}{\text{\Delta }t} = \dfrac{12000 - 11400}{2} = 300\text{m/s} \right.$.
}
\end{ex}

% ===== Câu 27 =====
% ID: L11C2B11-08-TLN
% Mức: TH
\dangbt{Nhận biết bản chất và tính chất của sóng điện từ}
\begin{ex}
% ID: L11C2B11-08-TLN
% Mức: TH
Ánh sáng đi $900\,\text{km}$ trong chân không. Thời gian tính theo ms bằng bao nhiêu?
\shortans{3}
\loigiai{
Xác định đúng dữ kiện và đơn vị, áp dụng hệ thức của dạng bài rồi tính được kết quả 3.
}
\end{ex}

% ===== Câu 28 =====
% ID: L11C2B11-08-TN
% Mức: VD
\dangbt{Bài toán truyền sóng điện từ qua các môi trường}
\begin{ex}
% ID: L11C2B11-08-TN
% Mức: VD
Trong chuyên đề “Bài toán truyền sóng điện từ qua các môi trường”, Kết luận nào phù hợp nhất khi giải bài toán sau: Ánh sáng đỏ có bước sóng 750 nm, ánh sáng tím 400 nm. Tính tần số mỗi ánh sáng và so sánh.
\choice
{Đại lượng cần tìm luôn bằng 0 trong mọi trường hợp.}
{Kết quả không phụ thuộc môi trường, tần số hoặc điều kiện biên.}
{\True $f_đ=4,0\times10^{14}$ Hz}
{Không cần sử dụng định nghĩa hay dữ kiện đã cho.}
\loigiai{
Đáp án C. $f_đ=4,0\times10^{14}$ Hz. Đối chiếu định nghĩa, hệ thức hoặc dữ kiện của bài toán cho kết luận trên.
}
\end{ex}

% ===== Câu 29 =====
% ID: L11C2B11-09-DS
% Mức: VD
\dangbt{Nhận biết bản chất và tính chất của sóng điện từ}
\begin{ex}
% ID: L11C2B11-09-DS
% Mức: VD
Xét các nhận định tổng hợp thuộc dạng “Nhận biết bản chất và tính chất của sóng điện từ”.
\choiceTF
{\True điện từ trường biến thiên.}
{\True truyền được với tốc độ $3\times10^8$ m/s.}
{\True vuông góc nhau và vuông góc phương truyền.}
{các phân tử mang điện đi xa.}
\loigiai{
\begin{itemchoice}
\item (Đúng) điện từ trường biến thiên. (Vì): iện từ trường biến thiên. b) Đ: truyền được với tốc độ $3\times10^8$ m/s. c) Đ: vuông góc nhau và vuông góc phương truyền. d) S: các phân tử mang điện đi xa
\item (Đúng) truyền được với tốc độ $3\times10^8$ m/s.
\item (Đúng) vuông góc nhau và vuông góc phương truyền.
\item (Sai) các phân tử mang điện đi xa.
\end{itemchoice}
}
\end{ex}

% ===== Câu 30 =====
% ID: L11C2B11-09-TL
% Mức: TH
\dangbt{Chưa có dạng}
\begin{ex}
% ID: L11C2B11-09-TL
% Mức: TH
*. Giả sử một vệ tinh truyền thông đang đứng yên so với mặt đất ở một độ cao xác định trong mặt phẳng xích đạo Trái Đất, đường thẳng nối vệ tinh với tâm Trái Đất đi qua kinh tuyến số 0 hoặc kinh tuyến gốc. Coi Trái Đất như một quả cầu bán kính $6400 \mathrm{km}$, khối lượng là $6.1024 \mathrm{kg}$ và chu kì quay quanh trục của nó là $24\,\mathrm{h}$, hằng số hấp dẫn $G = 6$,67 $\cdot$  10^{-11} $N\cdot$  m_2 / $\mathrm{kg}$ 2. Sóng cực ngắn $f \gt  30MHz phát từ vệ tinh truyền thẳng đến các điểm nằm trên xích đạo Trái Đất$ trong khoảng kinh độ nào? BÀl 12. GIAO THOA SÓNG $\nguon{ly11 kntt.pdf}$
\choiceTF

\loigiai{
Quỹ đạo của vệ tinh quanh Trái Đất được mô tả như Hình 11.2Ga.

[Giả sử một vệ tinh truyền thông đang đứng yên so với mặt đất]

Vì vệ tinh địa tĩnh đứng yên so với Trái Đất, lực hấp dẫn là lực hướng tâm, nên ta có: $\left. F_{hd} = F_{ht}\Leftrightarrow G\dfrac{Mm}{r^{2}} = m\left( \dfrac{2\pi}{T} \right)^{2}r \right.\left. \Rightarrow r = \sqrt[3]{GM\left( \dfrac{T}{2\pi} \right)^{2}} = \sqrt[3]{6,67 \cdot 10^{-11} \cdot 6 \cdot 10^{24}\left( \dfrac{24 \cdot 60 \cdot 60}{2\pi} \right)^{2}} \approx 42,3 \cdot 10^{6}\text{m}. \right.$

Vùng phủ sóng nằm trong miền giữa hai tiếp tuyến kẻ từ vệ tinh tới Trái Đất.

Do vậy, ta xác định được: $\left. \cos\varphi = \dfrac{\text{R}}{\text{r}} \approx \dfrac{1}{7}\Rightarrow\varphi \approx 81{^\circ}20^{\text{'}} \right.$ : Từ 81^(∘)20^(’) kinh độ tây đến kinh độ đông.
}
\end{ex}

% ===== Câu 31 =====
% ID: L11C2B11-09-TLN
% Mức: TH
\dangbt{Nhận biết bản chất và tính chất của sóng điện từ}
\begin{ex}
% ID: L11C2B11-09-TLN
% Mức: TH
Sóng điện từ có tốc độ $2\times10^8$ m/s trong môi trường. Trong 5 micro giây truyền được bao nhiêu m?
\shortans{1000}
\loigiai{
Xác định đúng dữ kiện và đơn vị, áp dụng hệ thức của dạng bài rồi tính được kết quả 1000.
}
\end{ex}

% ===== Câu 32 =====
% ID: L11C2B11-09-TN
% Mức: VD
\dangbt{Bài toán truyền sóng điện từ qua các môi trường}
\begin{ex}
% ID: L11C2B11-09-TN
% Mức: VD
Trong chuyên đề “Bài toán truyền sóng điện từ qua các môi trường”, Kết luận nào phù hợp nhất khi giải bài toán sau: Một xung radar đi tới mục tiêu rồi phản xạ về sau 200 micro giây. Tính khoảng cách tới mục tiêu.
\choice
{Không cần sử dụng định nghĩa hay dữ kiện đã cho.}
{Đại lượng cần tìm luôn bằng 0 trong mọi trường hợp.}
{Kết quả không phụ thuộc môi trường, tần số hoặc điều kiện biên.}
{\True Quãng đường hai chiều $ct$ nên $d=ct/2=3\times10^8\times200\times10^{-6}/2=30000$ m = $30\,\text{km}$.}
\loigiai{
Đáp án D. Quãng đường hai chiều $ct$ nên $d=ct/2=3\times10^8\times200\times10^{-6}/2=30000$ m = $30\,\text{km}$.. Đối chiếu định nghĩa, hệ thức hoặc dữ kiện của bài toán cho kết luận trên.
}
\end{ex}

% ===== Câu 33 =====
% ID: L11C2B11-10-DS
% Mức: VD
\dangbt{Nhận biết bản chất và tính chất của sóng điện từ}
\begin{ex}
% ID: L11C2B11-10-DS
% Mức: VD
Đánh giá cách xử lí các dữ kiện định lượng của dạng “Nhận biết bản chất và tính chất của sóng điện từ”.
\choiceTF
{\True Kết quả của bài toán thứ nhất là 600.}
{\True Kết quả của bài toán thứ hai là 3.}
{\True Kết quả của bài toán thứ ba là 1000.}
{Có thể bỏ qua hoàn toàn đơn vị và điều kiện vật lí khi tính toán.}
\loigiai{
\begin{itemchoice}
\item (Đúng) Kết quả của bài toán thứ nhất là 600. (Vì): Kết quả của bài toán thứ nhất là 600. b) Đ: Kết quả của bài toán thứ hai là 3. c) Đ: Kết quả của bài toán thứ ba là 1000. d) S: Có thể bỏ qua hoàn toàn đơn vị và điều kiện vật lí khi tính toán
\item (Đúng) Kết quả của bài toán thứ hai là 3.
\item (Đúng) Kết quả của bài toán thứ ba là 1000.
\item (Sai) Có thể bỏ qua hoàn toàn đơn vị và điều kiện vật lí khi tính toán.
\end{itemchoice}
}
\end{ex}

% ===== Câu 34 =====
% ID: L11C2B11-10-TL
% Mức: TH
\begin{ex}
% ID: L11C2B11-10-TL
% Mức: TH
Nêu ba điểm khác nhau cơ bản giữa sóng điện từ và sóng cơ.
\choiceTF

\loigiai{
Sóng điện từ là điện từ trường biến thiên, truyền được trong chân không và luôn là sóng ngang; sóng cơ là dao động vật chất, cần môi trường và có thể ngang hoặc dọc.
}
\end{ex}

% ===== Câu 35 =====
% ID: L11C2B11-10-TLN
% Mức: VD
\begin{ex}
% ID: L11C2B11-10-TLN
% Mức: VD
Từ kết quả của tình huống sau, nếu đại lượng cần tìm được nhân với hệ số 2 thì giá trị mới bằng bao nhiêu? Sóng điện từ truyền trong chân không $0,002\,\text{s}$. Quãng đường tính theo km bằng bao nhiêu?
\shortans{1200}
\loigiai{
Xác định đúng dữ kiện và đơn vị, áp dụng hệ thức của dạng bài rồi tính được kết quả 1200.
}
\end{ex}

% ===== Câu 36 =====
% ID: L11C2B11-10-TN
% Mức: VD
\dangbt{Bài toán truyền sóng điện từ qua các môi trường}
\begin{ex}
% ID: L11C2B11-10-TN
% Mức: VD
Trong chuyên đề “Bài toán truyền sóng điện từ qua các môi trường”, Phương pháp phù hợp nhất cho dạng “Bài toán truyền sóng điện từ qua các môi trường” là
\choice
{\True xác định dữ kiện, chọn đúng hệ thức hoặc định nghĩa, đổi đơn vị rồi kiểm tra kết quả}
{chỉ thay số mà không cần xét điều kiện áp dụng}
{bỏ qua đơn vị và chiều truyền sóng}
{luôn coi mọi điểm dao động cùng pha}
\loigiai{
Đáp án A. xác định dữ kiện, chọn đúng hệ thức hoặc định nghĩa, đổi đơn vị rồi kiểm tra kết quả. Đối chiếu định nghĩa, hệ thức hoặc dữ kiện của bài toán cho kết luận trên.
}
\end{ex}

% ===== Câu 37 =====
% ID: L11C2B11-11-DS
% Mức: TH
\dangbt{Nhận biết bức xạ điện từ theo nguồn phát và công dụng}
\begin{ex}
% ID: L11C2B11-11-DS
% Mức: TH
Trong chuyên đề “Nhận biết bức xạ điện từ theo nguồn phát và công dụng”, Sắp xếp thang sóng điện từ theo bước sóng giảm dần.
\choiceTF
{\True Sóng vô tuyến có bước sóng dài hơn ánh sáng nhìn thấy.}
{\True Tia hồng ngoại có bước sóng dài hơn tia tử ngoại.}
{\True Tia $X$ có bước sóng ngắn hơn tia tử ngoại.}
{Tia gamma có bước sóng dài nhất.}
\loigiai{
\begin{itemchoice}
\item (Đúng) Sóng vô tuyến có bước sóng dài hơn ánh sáng nhìn thấy. (Vì): Sóng vô tuyến có bước sóng dài hơn ánh sáng nhìn thấy. b) Đ: Tia hồng ngoại có bước sóng dài hơn tia tử ngoại. c) Đ: Tia $X$ có bước sóng ngắn hơn tia tử ngoại. d) S: Tia gamma có bước sóng dài nhất
\item (Đúng) Tia hồng ngoại có bước sóng dài hơn tia tử ngoại.
\item (Đúng) Tia $X$ có bước sóng ngắn hơn tia tử ngoại.
\item (Sai) Tia gamma có bước sóng dài nhất.
\end{itemchoice}
}
\end{ex}

% ===== Câu 38 =====
% ID: L11C2B11-11-TL
% Mức: TH
\dangbt{Nhận biết bản chất và tính chất của sóng điện từ}
\begin{bt}
% ID: L11C2B11-11-TL
% Mức: TH
Mô tả quan hệ phương và pha của $\vec E$, $\vec B$ và phương truyền sóng điện từ.
\loigiai{
$\vec E$ vuông góc $\vec B$; cả hai vuông góc phương truyền và biến thiên cùng pha tại một điểm.
}
\end{bt}

% ===== Câu 39 =====
% ID: L11C2B11-11-TLN
% Mức: VD
\begin{ex}
% ID: L11C2B11-11-TLN
% Mức: VD
Từ kết quả của tình huống sau, nếu đại lượng cần tìm được nhân với hệ số 0.5 thì giá trị mới bằng bao nhiêu? Ánh sáng đi $900\,\text{km}$ trong chân không. Thời gian tính theo ms bằng bao nhiêu?
\shortans{01/05/2026}
\loigiai{
Xác định đúng dữ kiện và đơn vị, áp dụng hệ thức của dạng bài rồi tính được kết quả 1.5.
}
\end{ex}

% ===== Câu 40 =====
% ID: L11C2B11-11-TN
% Mức: NB
\dangbt{Chưa có dạng}
\begin{ex}
% ID: L11C2B11-11-TN
% Mức: NB
Theo thứ tự bước sóng tăng dần thì sắp xếp nào dưới đây là đúng? $\nguon{ly11 kntt.pdf}$
\choice
{Vi sóng, tia tử ngoại, tia hồng ngoại, tia X.}
{\True Tia $X$, tia tử ngoại, tia hồng ngoại, vi sóng.}
{Tia tử ngoại, tia hồng ngoại, vi sóng, tia X.}
{Tia hồng ngoại, tia tử ngoại, vi sóng, tia X.}
\loigiai{
Chọn B.

Bước sóng tăng dần: Tia $X$, tia tử ngoại, tia hồng ngoại, vi sóng.
}
\end{ex}

% ===== Câu 41 =====
% ID: L11C2B11-12-DS
% Mức: TH
\dangbt{Nhận biết bức xạ điện từ theo nguồn phát và công dụng}
\begin{ex}
% ID: L11C2B11-12-DS
% Mức: TH
Trong chuyên đề “Nhận biết bức xạ điện từ theo nguồn phát và công dụng”, Về ứng dụng bức xạ điện từ.
\choiceTF
{\True Vi ba dùng trong thông tin và lò vi sóng.}
{\True Hồng ngoại dùng cảm biến nhiệt và điều khiển.}
{\True Tia $X$ dùng chẩn đoán hình ảnh.}
{Tử ngoại hoàn toàn vô hại với cơ thể.}
\loigiai{
\begin{itemchoice}
\item (Đúng) Vi ba dùng trong thông tin và lò vi sóng. (Vì): Vi ba dùng trong thông tin và lò vi sóng. b) Đ: Hồng ngoại dùng cảm biến nhiệt và điều khiển. c) Đ: Tia $X$ dùng chẩn đoán hình ảnh. d) S: Tử ngoại hoàn toàn vô hại với cơ thể
\item (Đúng) Hồng ngoại dùng cảm biến nhiệt và điều khiển.
\item (Đúng) Tia $X$ dùng chẩn đoán hình ảnh.
\item (Sai) Tử ngoại hoàn toàn vô hại với cơ thể.
\end{itemchoice}
}
\end{ex}

% ===== Câu 42 =====
% ID: L11C2B11-12-TL
% Mức: VD
\dangbt{Nhận biết bản chất và tính chất của sóng điện từ}
\begin{ex}
% ID: L11C2B11-12-TL
% Mức: VD
Giải thích vì sao ánh sáng được xem là sóng điện từ.
\choiceTF

\loigiai{
Ánh sáng có các tính chất sóng điện từ và truyền trong chân không với tốc độ đúng bằng tốc độ Maxwell dự đoán cho sóng điện từ.
}
\end{ex}

% ===== Câu 43 =====
% ID: L11C2B11-12-TLN
% Mức: VD
\begin{ex}
% ID: L11C2B11-12-TLN
% Mức: VD
Từ kết quả của tình huống sau, nếu đại lượng cần tìm được nhân với hệ số 1.5 thì giá trị mới bằng bao nhiêu? Sóng điện từ có tốc độ $2\times10^8$ m/s trong môi trường. Trong 5 micro giây truyền được bao nhiêu m?
\shortans{1500}
\loigiai{
Xác định đúng dữ kiện và đơn vị, áp dụng hệ thức của dạng bài rồi tính được kết quả 1500.
}
\end{ex}

% ===== Câu 44 =====
% ID: L11C2B11-12-TN
% Mức: NB
\dangbt{Chưa có dạng}
\begin{ex}
% ID: L11C2B11-12-TN
% Mức: NB
Phát biểu nào sau đây là sai khi nói về sóng điện từ? $\nguon{ly11 kntt.pdf}$
\choice
{Tất cả các sóng điện từ đều truyền trong chân không với tốc độ như nhau.}
{Sóng điện từ đều là sóng ngang.}
{Chúng đều tuân theo các quy luật phản xạ, khúc xạ.}
{\True Khi truyền từ không khí vào nước thì tần số, bước sóng và tốc độ của các sóng điện từ đều giảm.}
\loigiai{
Chọn D.

Khi truyền từ môi trường này sang môi trường kia tần số không đổi.
}
\end{ex}

% ===== Câu 45 =====
% ID: L11C2B11-13-DS
% Mức: TH
\dangbt{Nhận biết bức xạ điện từ theo nguồn phát và công dụng}
\begin{ex}
% ID: L11C2B11-13-DS
% Mức: TH
Trong chuyên đề “Nhận biết bức xạ điện từ theo nguồn phát và công dụng”, So sánh tần số và bước sóng trong chân không.
\choiceTF
{\True Tần số càng lớn thì bước sóng càng nhỏ.}
{\True Ánh sáng tím có tần số lớn hơn ánh sáng đỏ.}
{\True Tia gamma có năng lượng photon lớn hơn sóng vô tuyến.}
{Mọi bức xạ có cùng bước sóng.}
\loigiai{
\begin{itemchoice}
\item (Đúng) Tần số càng lớn thì bước sóng càng nhỏ. (Vì): Tần số càng lớn thì bước sóng càng nhỏ. b) Đ: Ánh sáng tím có tần số lớn hơn ánh sáng đỏ. c) Đ: Tia gamma có năng lượng photon lớn hơn sóng vô tuyến. d) S: Mọi bức xạ có cùng bước sóng
\item (Đúng) Ánh sáng tím có tần số lớn hơn ánh sáng đỏ.
\item (Đúng) Tia gamma có năng lượng photon lớn hơn sóng vô tuyến.
\item (Sai) Mọi bức xạ có cùng bước sóng.
\end{itemchoice}
}
\end{ex}

% ===== Câu 46 =====
% ID: L11C2B11-13-TL
% Mức: VD
\dangbt{Nhận biết bản chất và tính chất của sóng điện từ}
\begin{ex}
% ID: L11C2B11-13-TL
% Mức: VD
Phân tích các bước lập luận và chỉ ra điều kiện áp dụng trong bài toán: Nêu ba điểm khác nhau cơ bản giữa sóng điện từ và sóng cơ.
\choiceTF

\loigiai{
Sóng điện từ là điện từ trường biến thiên, truyền được trong chân không và luôn là sóng ngang; sóng cơ là dao động vật chất, cần môi trường và có thể ngang hoặc dọc. Cần nêu rõ điều kiện áp dụng, đổi đơn vị thống nhất và kiểm tra ý nghĩa vật lí của kết quả.
}
\end{ex}

% ===== Câu 47 =====
% ID: L11C2B11-13-TLN
% Mức: TH
\dangbt{Nhận biết bức xạ điện từ theo nguồn phát và công dụng}
\begin{ex}
% ID: L11C2B11-13-TLN
% Mức: TH
Trong chuyên đề “Nhận biết bức xạ điện từ theo nguồn phát và công dụng”, Bức xạ có bước sóng 600 nm. Đổi sang micromet bằng bao nhiêu?
\shortans{0.6}
\loigiai{
Xác định đúng dữ kiện và đơn vị, áp dụng hệ thức của dạng bài rồi tính được kết quả 0.6.
}
\end{ex}

% ===== Câu 48 =====
% ID: L11C2B11-13-TN
% Mức: NB
\dangbt{Chưa có dạng}
\begin{ex}
% ID: L11C2B11-13-TN
% Mức: NB
Sóng điện từ có bước sóng nào dưới đây thuộc về tia hồng ngoại? $\nguon{ly11 kntt.pdf}$
\choice
{$7\cdot 10^{-2} m.$}
{\True $7\cdot 10^{-6} m.$}
{$7\cdot 10^{-9} m.$}
{$7 \cdot  10^{-12} m.$}
\loigiai{
Chọn B.

Tia hồng ngoại có bước sóng lớn hơn bước sóng của ánh sáng đỏ.
}
\end{ex}

% ===== Câu 49 =====
% ID: L11C2B11-14-DS
% Mức: VD
\dangbt{Nhận biết bức xạ điện từ theo nguồn phát và công dụng}
\begin{ex}
% ID: L11C2B11-14-DS
% Mức: VD
Trong chuyên đề “Nhận biết bức xạ điện từ theo nguồn phát và công dụng”, Xét các nhận định tổng hợp thuộc dạng “Nhận biết bức xạ điện từ theo nguồn phát và công dụng”.
\choiceTF
{\True sóng vô tuyến.}
{\True tia hồng ngoại.}
{\True chụp ảnh xương.}
{tia tử ngoại.}
\loigiai{
\begin{itemchoice}
\item (Đúng) sóng vô tuyến. (Vì): sóng vô tuyến. b) Đ: tia hồng ngoại. c) Đ: chụp ảnh xương. d) S: tia tử ngoại
\item (Đúng) tia hồng ngoại.
\item (Đúng) chụp ảnh xương.
\item (Sai) tia tử ngoại.
\end{itemchoice}
}
\end{ex}

% ===== Câu 50 =====
% ID: L11C2B11-14-TL
% Mức: VD
\dangbt{Nhận biết bản chất và tính chất của sóng điện từ}
\begin{ex}
% ID: L11C2B11-14-TL
% Mức: VD
Một học sinh giải bài sau nhưng bỏ qua điều kiện vật lí. Hãy sửa cách làm và trình bày lời giải đúng: Mô tả quan hệ phương và pha của $\vec E$, $\vec B$ và phương truyền sóng điện từ.
\choiceTF

\loigiai{
Cách giải đúng: $\vec E$ vuông góc $\vec B$; cả hai vuông góc phương truyền và biến thiên cùng pha tại một điểm. Sai lầm cần tránh là dùng hệ thức khi chưa xác định điều kiện biên, môi trường hoặc đơn vị.
}
\end{ex}

% ===== Câu 51 =====
% ID: L11C2B11-14-TLN
% Mức: TH
\dangbt{Nhận biết bức xạ điện từ theo nguồn phát và công dụng}
\begin{ex}
% ID: L11C2B11-14-TLN
% Mức: TH
Trong chuyên đề “Nhận biết bức xạ điện từ theo nguồn phát và công dụng”, Tia có tần số $6\times10^{14}$ Hz trong chân không. Bước sóng tính theo nm bằng bao nhiêu?
\shortans{500}
\loigiai{
Xác định đúng dữ kiện và đơn vị, áp dụng hệ thức của dạng bài rồi tính được kết quả 500.
}
\end{ex}

% ===== Câu 52 =====
% ID: L11C2B11-14-TN
% Mức: NB
\dangbt{Chưa có dạng}
\begin{ex}
% ID: L11C2B11-14-TN
% Mức: NB
Một sóng vô tuyến có tần số $108 \mathrm{Hz}$ được truyền trong không trung với tốc độ $3.108 \mathrm{m/s}$. Bước sóng của sóng đó là $\nguon{ly11 kntt.pdf}$
\choice
{$1,5\,\text{m}$.}
{\True $3\,\text{m}$.}
{$0,33\,\text{m}$.}
{$0,16\,\text{m}$. 8}
\loigiai{
Chọn B.

Bước sóng $\lambda = \dfrac{3\cdot 10^{8}}{10^{8}} =$ 3 $\mathrm{m}$
}
\end{ex}

% ===== Câu 53 =====
% ID: L11C2B11-15-DS
% Mức: VD
\dangbt{Nhận biết bức xạ điện từ theo nguồn phát và công dụng}
\begin{ex}
% ID: L11C2B11-15-DS
% Mức: VD
Trong chuyên đề “Nhận biết bức xạ điện từ theo nguồn phát và công dụng”, Đánh giá cách xử lí các dữ kiện định lượng của dạng “Nhận biết bức xạ điện từ theo nguồn phát và công dụng”.
\choiceTF
{\True Kết quả của bài toán thứ nhất là 0.6.}
{\True Kết quả của bài toán thứ hai là 500.}
{\True Kết quả của bài toán thứ ba là 3.}
{Có thể bỏ qua hoàn toàn đơn vị và điều kiện vật lí khi tính toán.}
\loigiai{
\begin{itemchoice}
\item (Đúng) Kết quả của bài toán thứ nhất là 0.6. (Vì): Kết quả của bài toán thứ nhất là 0.6. b) Đ: Kết quả của bài toán thứ hai là 500. c) Đ: Kết quả của bài toán thứ ba là 3. d) S: Có thể bỏ qua hoàn toàn đơn vị và điều kiện vật lí khi tính toán
\item (Đúng) Kết quả của bài toán thứ hai là 500.
\item (Đúng) Kết quả của bài toán thứ ba là 3.
\item (Sai) Có thể bỏ qua hoàn toàn đơn vị và điều kiện vật lí khi tính toán.
\end{itemchoice}
}
\end{ex}

% ===== Câu 54 =====
% ID: L11C2B11-15-TL
% Mức: VD
\dangbt{Nhận biết bản chất và tính chất của sóng điện từ}
\begin{ex}
% ID: L11C2B11-15-TL
% Mức: VD
Mở rộng tình huống sau thành một ví dụ thực tế và trình bày phương pháp giải: Giải thích vì sao ánh sáng được xem là sóng điện từ.
\choiceTF

\loigiai{
Ánh sáng có các tính chất sóng điện từ và truyền trong chân không với tốc độ đúng bằng tốc độ Maxwell dự đoán cho sóng điện từ. Có thể kiểm chứng bằng cách thay kết quả vào dữ kiện ban đầu và đối chiếu với hiện tượng thực tế.
}
\end{ex}

% ===== Câu 55 =====
% ID: L11C2B11-15-TLN
% Mức: TH
\dangbt{Nhận biết bức xạ điện từ theo nguồn phát và công dụng}
\begin{ex}
% ID: L11C2B11-15-TLN
% Mức: TH
Trong chuyên đề “Nhận biết bức xạ điện từ theo nguồn phát và công dụng”, Sóng vô tuyến có bước sóng $100\,\text{m}$. Tần số tính theo MHz bằng bao nhiêu?
\shortans{3}
\loigiai{
Xác định đúng dữ kiện và đơn vị, áp dụng hệ thức của dạng bài rồi tính được kết quả 3.
}
\end{ex}

% ===== Câu 56 =====
% ID: L11C2B11-15-TN
% Mức: NB
\dangbt{Chưa có dạng}
\begin{ex}
% ID: L11C2B11-15-TN
% Mức: NB
Sóng vô tuyến truyền trong không trung với tốc độ 3 $\cdot10 \mathrm{m/s}$. Một đài phát sóng radio có tần số $106 \mathrm{Hz}$. Bước sóng của sóng radio này là $\nguon{ly11 kntt.pdf}$
\choice
{\True $300\,\text{m}$.}
{$150\,\text{m}$.}
{$0,30\,\text{m}$.}
{$0,15\,\text{m}$}
\loigiai{
Chọn A.

Bước sóng $\lambda = \dfrac{3\cdot 10^{8}}{10^{6}} =$ 300 $\mathrm{m}$
}
\end{ex}

% ===== Câu 57 =====
% ID: L11C2B11-16-DS
% Mức: TH
\dangbt{Phân loại thang sóng điện từ và ứng dụng}
\begin{ex}
% ID: L11C2B11-16-DS
% Mức: TH
Sắp xếp thang sóng điện từ theo bước sóng giảm dần.
\choiceTF
{\True Sóng vô tuyến có bước sóng dài hơn ánh sáng nhìn thấy.}
{\True Tia hồng ngoại có bước sóng dài hơn tia tử ngoại.}
{\True Tia $X$ có bước sóng ngắn hơn tia tử ngoại.}
{Tia gamma có bước sóng dài nhất.}
\loigiai{
\begin{itemchoice}
\item (Đúng) Sóng vô tuyến có bước sóng dài hơn ánh sáng nhìn thấy. (Vì): Sóng vô tuyến có bước sóng dài hơn ánh sáng nhìn thấy. b) Đ: Tia hồng ngoại có bước sóng dài hơn tia tử ngoại. c) Đ: Tia $X$ có bước sóng ngắn hơn tia tử ngoại. d) S: Tia gamma có bước sóng dài nhất
\item (Đúng) Tia hồng ngoại có bước sóng dài hơn tia tử ngoại.
\item (Đúng) Tia $X$ có bước sóng ngắn hơn tia tử ngoại.
\item (Sai) Tia gamma có bước sóng dài nhất.
\end{itemchoice}
}
\end{ex}

% ===== Câu 58 =====
% ID: L11C2B11-16-TL
% Mức: TH
\dangbt{Nhận biết bức xạ điện từ theo nguồn phát và công dụng}
\begin{ex}
% ID: L11C2B11-16-TL
% Mức: TH
Trong chuyên đề “Nhận biết bức xạ điện từ theo nguồn phát và công dụng”, Lập thứ tự từ bước sóng dài đến ngắn cho: hồng ngoại, tử ngoại, sóng vô tuyến, ánh sáng nhìn thấy, tia X.
\choiceTF

\loigiai{
Sóng vô tuyến, hồng ngoại, ánh sáng nhìn thấy, tử ngoại, tia X.
}
\end{ex}

% ===== Câu 59 =====
% ID: L11C2B11-16-TLN
% Mức: VD
\begin{ex}
% ID: L11C2B11-16-TLN
% Mức: VD
Trong chuyên đề “Nhận biết bức xạ điện từ theo nguồn phát và công dụng”, Từ kết quả của tình huống sau, nếu đại lượng cần tìm được nhân với hệ số 2 thì giá trị mới bằng bao nhiêu? Bức xạ có bước sóng 600 nm. Đổi sang micromet bằng bao nhiêu?
\shortans{01/02/2026}
\loigiai{
Xác định đúng dữ kiện và đơn vị, áp dụng hệ thức của dạng bài rồi tính được kết quả 1.2.
}
\end{ex}

% ===== Câu 60 =====
% ID: L11C2B11-16-TN
% Mức: TH
\dangbt{Chưa có dạng}
\begin{ex}
% ID: L11C2B11-16-TN
% Mức: TH
Nội dung nào sau đây tóm tắt đúng đặc điểm của sóng điện từ, tính từ sóng vô tuyến đến tia $\gamma$ trong thang của sóng điện từ? Tần số Bước sóng Tốc độ trong chân không $\nguon{ly11 kntt.pdf}$
\choice
{tăng dần giảm dần giảm dần}
{giảm dần tăng dần tăng dần}
{\True tăng dần giảm dần không đổi}
{giảm dần tăng dần không đổi}
\loigiai{
Chọn C.

Tính từ sóng vô tuyến đến tia γ thì bước sóng giảm dần, tần số tăng dần, tốc độ trong chân không là như nhau.
}
\end{ex}

% ===== Câu 61 =====
% ID: L11C2B11-17-DS
% Mức: TH
\dangbt{Phân loại thang sóng điện từ và ứng dụng}
\begin{ex}
% ID: L11C2B11-17-DS
% Mức: TH
Về ứng dụng bức xạ điện từ.
\choiceTF
{\True Vi ba dùng trong thông tin và lò vi sóng.}
{\True Hồng ngoại dùng cảm biến nhiệt và điều khiển.}
{\True Tia $X$ dùng chẩn đoán hình ảnh.}
{Tử ngoại hoàn toàn vô hại với cơ thể.}
\loigiai{
\begin{itemchoice}
\item (Đúng) Vi ba dùng trong thông tin và lò vi sóng. (Vì): Vi ba dùng trong thông tin và lò vi sóng. b) Đ: Hồng ngoại dùng cảm biến nhiệt và điều khiển. c) Đ: Tia $X$ dùng chẩn đoán hình ảnh. d) S: Tử ngoại hoàn toàn vô hại với cơ thể
\item (Đúng) Hồng ngoại dùng cảm biến nhiệt và điều khiển.
\item (Đúng) Tia $X$ dùng chẩn đoán hình ảnh.
\item (Sai) Tử ngoại hoàn toàn vô hại với cơ thể.
\end{itemchoice}
}
\end{ex}

% ===== Câu 62 =====
% ID: L11C2B11-17-TL
% Mức: TH
\dangbt{Nhận biết bức xạ điện từ theo nguồn phát và công dụng}
\begin{ex}
% ID: L11C2B11-17-TL
% Mức: TH
Trong chuyên đề “Nhận biết bức xạ điện từ theo nguồn phát và công dụng”, Chọn bức xạ phù hợp cho điều khiển từ xa, chụp xương và truyền thông vệ tinh; giải thích ngắn gọn.
\choiceTF

\loigiai{
Điều khiển từ xa: hồng ngoại; chụp xương: tia $X$ do khả năng đâm xuyên khác nhau; vệ tinh: vi ba do truyền định hướng và qua khí quyển phù hợp.
}
\end{ex}

% ===== Câu 63 =====
% ID: L11C2B11-17-TLN
% Mức: VD
\begin{ex}
% ID: L11C2B11-17-TLN
% Mức: VD
Trong chuyên đề “Nhận biết bức xạ điện từ theo nguồn phát và công dụng”, Từ kết quả của tình huống sau, nếu đại lượng cần tìm được nhân với hệ số 0.5 thì giá trị mới bằng bao nhiêu? Tia có tần số $6\times10^{14}$ Hz trong chân không. Bước sóng tính theo nm bằng bao nhiêu?
\shortans{250}
\loigiai{
Xác định đúng dữ kiện và đơn vị, áp dụng hệ thức của dạng bài rồi tính được kết quả 250.
}
\end{ex}

% ===== Câu 64 =====
% ID: L11C2B11-17-TN
% Mức: TH
\dangbt{Chưa có dạng}
\begin{ex}
% ID: L11C2B11-17-TN
% Mức: TH
Một sóng ánh sáng có bước sóng $\lambda$ _1 và tốc độ v_1 khi truyền trong chân không. Khi đi vào trong tấm thuỷ tinh có bước sóng $\lambda$ _2 và tốc độ v_2. Biểu thức nào dưới đây biểu diễn đúng mối liên hệ giữa v_2 với $\lambda$ _1, $\lambda$ _2 và v_1 ? $\lambda\lambda\lambda\lambda\nguon{ly11 kntt.pdf}$
\choice
{$v_2 = \lambda_1 \cdot  v_1$.}
{\True $v_2 = \lambda_2 \cdot  v_1$.}
{$v_2 = 2v 1$.}
{$v_2 = \lambda_2 \lambda_1 v_1$. 2 1 1}
\loigiai{
Chọn B.

Khi truyền từ môi trường này sang môi trường kia tần số không đổi.

$\left. \dfrac{v_{1}}{\lambda_{1}} = \dfrac{v_{2}}{\lambda_{2}}\Rightarrow v_{2} = v_{1}\dfrac{\lambda_{2}}{\lambda_{1}} \right.$
}
\end{ex}

% ===== Câu 65 =====
% ID: L11C2B11-18-DS
% Mức: TH
\dangbt{Phân loại thang sóng điện từ và ứng dụng}
\begin{ex}
% ID: L11C2B11-18-DS
% Mức: TH
So sánh tần số và bước sóng trong chân không.
\choiceTF
{\True Tần số càng lớn thì bước sóng càng nhỏ.}
{\True Ánh sáng tím có tần số lớn hơn ánh sáng đỏ.}
{\True Tia gamma có năng lượng photon lớn hơn sóng vô tuyến.}
{Mọi bức xạ có cùng bước sóng.}
\loigiai{
\begin{itemchoice}
\item (Đúng) Tần số càng lớn thì bước sóng càng nhỏ. (Vì): Tần số càng lớn thì bước sóng càng nhỏ. b) Đ: Ánh sáng tím có tần số lớn hơn ánh sáng đỏ. c) Đ: Tia gamma có năng lượng photon lớn hơn sóng vô tuyến. d) S: Mọi bức xạ có cùng bước sóng
\item (Đúng) Ánh sáng tím có tần số lớn hơn ánh sáng đỏ.
\item (Đúng) Tia gamma có năng lượng photon lớn hơn sóng vô tuyến.
\item (Sai) Mọi bức xạ có cùng bước sóng.
\end{itemchoice}
}
\end{ex}

% ===== Câu 66 =====
% ID: L11C2B11-18-TL
% Mức: VD
\dangbt{Nhận biết bức xạ điện từ theo nguồn phát và công dụng}
\begin{ex}
% ID: L11C2B11-18-TL
% Mức: VD
Trong chuyên đề “Nhận biết bức xạ điện từ theo nguồn phát và công dụng”, Giải thích vì sao cần hạn chế tiếp xúc tia tử ngoại và tia $X$ dù chúng có nhiều ứng dụng.
\choiceTF

\loigiai{
Đây là các bức xạ tần số cao, có thể gây tổn thương mô và vật liệu di truyền; cần kiểm soát liều và che chắn.
}
\end{ex}

% ===== Câu 67 =====
% ID: L11C2B11-18-TLN
% Mức: VD
\begin{ex}
% ID: L11C2B11-18-TLN
% Mức: VD
Trong chuyên đề “Nhận biết bức xạ điện từ theo nguồn phát và công dụng”, Từ kết quả của tình huống sau, nếu đại lượng cần tìm được nhân với hệ số 1.5 thì giá trị mới bằng bao nhiêu? Sóng vô tuyến có bước sóng $100\,\text{m}$. Tần số tính theo MHz bằng bao nhiêu?
\shortans{04/05/2026}
\loigiai{
Xác định đúng dữ kiện và đơn vị, áp dụng hệ thức của dạng bài rồi tính được kết quả 4.5.
}
\end{ex}

% ===== Câu 68 =====
% ID: L11C2B11-18-TN
% Mức: NB
\dangbt{Nhận biết bản chất và tính chất của sóng điện từ}
\begin{ex}
% ID: L11C2B11-18-TN
% Mức: NB
Sóng điện từ là sự lan truyền của
\choice
{áp suất cơ học}
{dòng điện không đổi}
{\True điện từ trường biến thiên}
{các phân tử mang điện đi xa}
\loigiai{
Đáp án C. điện từ trường biến thiên. Đối chiếu định nghĩa, hệ thức hoặc dữ kiện của bài toán cho kết luận trên.
}
\end{ex}

% ===== Câu 69 =====
% ID: L11C2B11-19-DS
% Mức: VD
\dangbt{Phân loại thang sóng điện từ và ứng dụng}
\begin{ex}
% ID: L11C2B11-19-DS
% Mức: VD
Xét các nhận định tổng hợp thuộc dạng “Phân loại thang sóng điện từ và ứng dụng”.
\choiceTF
{\True sóng vô tuyến.}
{\True tia hồng ngoại.}
{\True chụp ảnh xương.}
{tia tử ngoại.}
\loigiai{
\begin{itemchoice}
\item (Đúng) sóng vô tuyến. (Vì): sóng vô tuyến. b) Đ: tia hồng ngoại. c) Đ: chụp ảnh xương. d) S: tia tử ngoại
\item (Đúng) tia hồng ngoại.
\item (Đúng) chụp ảnh xương.
\item (Sai) tia tử ngoại.
\end{itemchoice}
}
\end{ex}

% ===== Câu 70 =====
% ID: L11C2B11-19-TL
% Mức: VD
\dangbt{Nhận biết bức xạ điện từ theo nguồn phát và công dụng}
\begin{ex}
% ID: L11C2B11-19-TL
% Mức: VD
Trong chuyên đề “Nhận biết bức xạ điện từ theo nguồn phát và công dụng”, Phân tích các bước lập luận và chỉ ra điều kiện áp dụng trong bài toán: Lập thứ tự từ bước sóng dài đến ngắn cho: hồng ngoại, tử ngoại, sóng vô tuyến, ánh sáng nhìn thấy, tia X.
\choiceTF

\loigiai{
Sóng vô tuyến, hồng ngoại, ánh sáng nhìn thấy, tử ngoại, tia X. Cần nêu rõ điều kiện áp dụng, đổi đơn vị thống nhất và kiểm tra ý nghĩa vật lí của kết quả.
}
\end{ex}

% ===== Câu 71 =====
% ID: L11C2B11-19-TLN
% Mức: TH
\dangbt{Phân loại thang sóng điện từ và ứng dụng}
\begin{ex}
% ID: L11C2B11-19-TLN
% Mức: TH
Bức xạ có bước sóng 600 nm. Đổi sang micromet bằng bao nhiêu?
\shortans{0.6}
\loigiai{
Xác định đúng dữ kiện và đơn vị, áp dụng hệ thức của dạng bài rồi tính được kết quả 0.6.
}
\end{ex}

% ===== Câu 72 =====
% ID: L11C2B11-19-TN
% Mức: NB
\dangbt{Nhận biết bản chất và tính chất của sóng điện từ}
\begin{ex}
% ID: L11C2B11-19-TN
% Mức: NB
Trong chân không, sóng điện từ
\choice
{không truyền được}
{chỉ là sóng dọc}
{có tốc độ phụ thuộc tần số rất mạnh}
{\True truyền được với tốc độ $3\times10^8$ m/s}
\loigiai{
Đáp án D. truyền được với tốc độ $3\times10^8$ m/s. Đối chiếu định nghĩa, hệ thức hoặc dữ kiện của bài toán cho kết luận trên.
}
\end{ex}

% ===== Câu 73 =====
% ID: L11C2B11-20-DS
% Mức: VD
\dangbt{Phân loại thang sóng điện từ và ứng dụng}
\begin{ex}
% ID: L11C2B11-20-DS
% Mức: VD
Đánh giá cách xử lí các dữ kiện định lượng của dạng “Phân loại thang sóng điện từ và ứng dụng”.
\choiceTF
{\True Kết quả của bài toán thứ nhất là 0.6.}
{\True Kết quả của bài toán thứ hai là 500.}
{\True Kết quả của bài toán thứ ba là 3.}
{Có thể bỏ qua hoàn toàn đơn vị và điều kiện vật lí khi tính toán.}
\loigiai{
\begin{itemchoice}
\item (Đúng) Kết quả của bài toán thứ nhất là 0.6. (Vì): Kết quả của bài toán thứ nhất là 0.6. b) Đ: Kết quả của bài toán thứ hai là 500. c) Đ: Kết quả của bài toán thứ ba là 3. d) S: Có thể bỏ qua hoàn toàn đơn vị và điều kiện vật lí khi tính toán
\item (Đúng) Kết quả của bài toán thứ hai là 500.
\item (Đúng) Kết quả của bài toán thứ ba là 3.
\item (Sai) Có thể bỏ qua hoàn toàn đơn vị và điều kiện vật lí khi tính toán.
\end{itemchoice}
}
\end{ex}

% ===== Câu 74 =====
% ID: L11C2B11-20-TL
% Mức: VD
\dangbt{Nhận biết bức xạ điện từ theo nguồn phát và công dụng}
\begin{ex}
% ID: L11C2B11-20-TL
% Mức: VD
Trong chuyên đề “Nhận biết bức xạ điện từ theo nguồn phát và công dụng”, Một học sinh giải bài sau nhưng bỏ qua điều kiện vật lí. Hãy sửa cách làm và trình bày lời giải đúng: Chọn bức xạ phù hợp cho điều khiển từ xa, chụp xương và truyền thông vệ tinh; giải thích ngắn gọn.
\choiceTF

\loigiai{
Cách giải đúng: Điều khiển từ xa: hồng ngoại; chụp xương: tia $X$ do khả năng đâm xuyên khác nhau; vệ tinh: vi ba do truyền định hướng và qua khí quyển phù hợp. Sai lầm cần tránh là dùng hệ thức khi chưa xác định điều kiện biên, môi trường hoặc đơn vị.
}
\end{ex}

% ===== Câu 75 =====
% ID: L11C2B11-20-TLN
% Mức: TH
\dangbt{Phân loại thang sóng điện từ và ứng dụng}
\begin{ex}
% ID: L11C2B11-20-TLN
% Mức: TH
Tia có tần số $6\times10^{14}$ Hz trong chân không. Bước sóng tính theo nm bằng bao nhiêu?
\shortans{500}
\loigiai{
Xác định đúng dữ kiện và đơn vị, áp dụng hệ thức của dạng bài rồi tính được kết quả 500.
}
\end{ex}

% ===== Câu 76 =====
% ID: L11C2B11-20-TN
% Mức: TH
\dangbt{Nhận biết bản chất và tính chất của sóng điện từ}
\begin{ex}
% ID: L11C2B11-20-TN
% Mức: TH
Trong sóng điện từ, vectơ điện trường và vectơ từ trường
\choice
{\True vuông góc nhau và vuông góc phương truyền}
{song song nhau}
{cùng phương truyền}
{luôn bằng 0}
\loigiai{
Đáp án A. vuông góc nhau và vuông góc phương truyền. Đối chiếu định nghĩa, hệ thức hoặc dữ kiện của bài toán cho kết luận trên.
}
\end{ex}

% ===== Câu 77 =====
% ID: L11C2B11-21-DS
% Mức: TH
\dangbt{Tính bước sóng, tần số và tốc độ của sóng điện từ}
\begin{ex}
% ID: L11C2B11-21-DS
% Mức: TH
Sóng điện từ có $f=150$ MHz trong chân không.
\choiceTF
{\True Bước sóng bằng $2\,\text{m}$.}
{\True Chu kì bằng $6,67$ ns.}
{\True Trong 1 micro giây sóng truyền $300\,\text{m}$.}
{Đổi sang môi trường khác làm tần số giảm.}
\loigiai{
\begin{itemchoice}
\item (Đúng) Bước sóng bằng $2\,\text{m}$. (Vì): Bước sóng bằng $2\,\text{m}$. b) Đ: Chu kì bằng $6,67$ ns. c) Đ: Trong 1 micro giây sóng truyền $300\,\text{m}$. d) S: Đổi sang môi trường khác làm tần số giảm
\item (Đúng) Chu kì bằng $6,67$ ns.
\item (Đúng) Trong 1 micro giây sóng truyền $300\,\text{m}$.
\item (Sai) Đổi sang môi trường khác làm tần số giảm.
\end{itemchoice}
}
\end{ex}

% ===== Câu 78 =====
% ID: L11C2B11-21-TL
% Mức: VD
\dangbt{Nhận biết bức xạ điện từ theo nguồn phát và công dụng}
\begin{ex}
% ID: L11C2B11-21-TL
% Mức: VD
Trong chuyên đề “Nhận biết bức xạ điện từ theo nguồn phát và công dụng”, Mở rộng tình huống sau thành một ví dụ thực tế và trình bày phương pháp giải: Giải thích vì sao cần hạn chế tiếp xúc tia tử ngoại và tia $X$ dù chúng có nhiều ứng dụng.
\choiceTF

\loigiai{
Đây là các bức xạ tần số cao, có thể gây tổn thương mô và vật liệu di truyền; cần kiểm soát liều và che chắn. Có thể kiểm chứng bằng cách thay kết quả vào dữ kiện ban đầu và đối chiếu với hiện tượng thực tế.
}
\end{ex}

% ===== Câu 79 =====
% ID: L11C2B11-21-TLN
% Mức: TH
\dangbt{Phân loại thang sóng điện từ và ứng dụng}
\begin{ex}
% ID: L11C2B11-21-TLN
% Mức: TH
Sóng vô tuyến có bước sóng $100\,\text{m}$. Tần số tính theo MHz bằng bao nhiêu?
\shortans{3}
\loigiai{
Xác định đúng dữ kiện và đơn vị, áp dụng hệ thức của dạng bài rồi tính được kết quả 3.
}
\end{ex}

% ===== Câu 80 =====
% ID: L11C2B11-21-TN
% Mức: TH
\dangbt{Nhận biết bản chất và tính chất của sóng điện từ}
\begin{ex}
% ID: L11C2B11-21-TN
% Mức: TH
Sóng điện từ truyền trong chân không $0,002\,\text{s}$. Quãng đường tính theo km bằng bao nhiêu? Chọn giá trị đúng.
\choice
{601}
{\True 600}
{1200}
{300}
\loigiai{
Đáp án B. 600. Đối chiếu định nghĩa, hệ thức hoặc dữ kiện của bài toán cho kết luận trên.
}
\end{ex}

% ===== Câu 81 =====
% ID: L11C2B11-22-DS
% Mức: TH
\dangbt{Tính bước sóng, tần số và tốc độ của sóng điện từ}
\begin{ex}
% ID: L11C2B11-22-DS
% Mức: TH
Ánh sáng có bước sóng 500 nm trong chân không.
\choiceTF
{\True Tần số bằng $6\times10^{14}$ Hz.}
{\True Đây thuộc vùng ánh sáng nhìn thấy.}
{\True Chu kì xấp xỉ $1,67\times10^{-15}$ s.}
{Trong điện môi có tốc độ nhỏ hơn thì tần số nhỏ hơn.}
\loigiai{
\begin{itemchoice}
\item (Đúng) Tần số bằng $6\times10^{14}$ Hz. (Vì): Tần số bằng $6\times10^{14}$ Hz. b) Đ: Đây thuộc vùng ánh sáng nhìn thấy. c) Đ: Chu kì xấp xỉ $1,67\times10^{-15}$ s. d) S: Trong điện môi có tốc độ nhỏ hơn thì tần số nhỏ hơn
\item (Đúng) Đây thuộc vùng ánh sáng nhìn thấy.
\item (Đúng) Chu kì xấp xỉ $1,67\times10^{-15}$ s.
\item (Sai) Trong điện môi có tốc độ nhỏ hơn thì tần số nhỏ hơn.
\end{itemchoice}
}
\end{ex}

% ===== Câu 82 =====
% ID: L11C2B11-22-TL
% Mức: TH
\dangbt{Phân loại thang sóng điện từ và ứng dụng}
\begin{ex}
% ID: L11C2B11-22-TL
% Mức: TH
Lập thứ tự từ bước sóng dài đến ngắn cho: hồng ngoại, tử ngoại, sóng vô tuyến, ánh sáng nhìn thấy, tia X.
\choiceTF

\loigiai{
Sóng vô tuyến, hồng ngoại, ánh sáng nhìn thấy, tử ngoại, tia X.
}
\end{ex}

% ===== Câu 83 =====
% ID: L11C2B11-22-TLN
% Mức: VD
\begin{ex}
% ID: L11C2B11-22-TLN
% Mức: VD
Từ kết quả của tình huống sau, nếu đại lượng cần tìm được nhân với hệ số 2 thì giá trị mới bằng bao nhiêu? Bức xạ có bước sóng 600 nm. Đổi sang micromet bằng bao nhiêu?
\shortans{01/02/2026}
\loigiai{
Xác định đúng dữ kiện và đơn vị, áp dụng hệ thức của dạng bài rồi tính được kết quả 1.2.
}
\end{ex}

% ===== Câu 84 =====
% ID: L11C2B11-22-TN
% Mức: TH
\dangbt{Nhận biết bản chất và tính chất của sóng điện từ}
\begin{ex}
% ID: L11C2B11-22-TN
% Mức: TH
Ánh sáng đi $900\,\text{km}$ trong chân không. Thời gian tính theo ms bằng bao nhiêu? Chọn giá trị đúng.
\choice
{01/05/2026}
{4}
{\True 3}
{6}
\loigiai{
Đáp án C. 3. Đối chiếu định nghĩa, hệ thức hoặc dữ kiện của bài toán cho kết luận trên.
}
\end{ex}

% ===== Câu 85 =====
% ID: L11C2B11-23-DS
% Mức: TH
\dangbt{Tính bước sóng, tần số và tốc độ của sóng điện từ}
\begin{ex}
% ID: L11C2B11-23-DS
% Mức: TH
Sóng 75 MHz truyền trong môi trường với tốc độ $2,25\times10^8$ m/s.
\choiceTF
{\True Bước sóng bằng $3\,\text{m}$.}
{\True Chu kì bằng $1,33\times10^{-8}$ s.}
{\True Tần số không đổi khi truyền sang chân không.}
{Bước sóng trong chân không vẫn bằng $3\,\text{m}$.}
\loigiai{
\begin{itemchoice}
\item (Đúng) Bước sóng bằng $3\,\text{m}$. (Vì): Bước sóng bằng $3\,\text{m}$. b) Đ: Chu kì bằng $1,33\times10^{-8}$ s. c) Đ: Tần số không đổi khi truyền sang chân không. d) S: Bước sóng trong chân không vẫn bằng $3\,\text{m}$
\item (Đúng) Chu kì bằng $1,33\times10^{-8}$ s.
\item (Đúng) Tần số không đổi khi truyền sang chân không.
\item (Sai) Bước sóng trong chân không vẫn bằng $3\,\text{m}$.
\end{itemchoice}
}
\end{ex}

% ===== Câu 86 =====
% ID: L11C2B11-23-TL
% Mức: TH
\dangbt{Phân loại thang sóng điện từ và ứng dụng}
\begin{ex}
% ID: L11C2B11-23-TL
% Mức: TH
Chọn bức xạ phù hợp cho điều khiển từ xa, chụp xương và truyền thông vệ tinh; giải thích ngắn gọn.
\choiceTF

\loigiai{
Điều khiển từ xa: hồng ngoại; chụp xương: tia $X$ do khả năng đâm xuyên khác nhau; vệ tinh: vi ba do truyền định hướng và qua khí quyển phù hợp.
}
\end{ex}

% ===== Câu 87 =====
% ID: L11C2B11-23-TLN
% Mức: VD
\begin{ex}
% ID: L11C2B11-23-TLN
% Mức: VD
Từ kết quả của tình huống sau, nếu đại lượng cần tìm được nhân với hệ số 0.5 thì giá trị mới bằng bao nhiêu? Tia có tần số $6\times10^{14}$ Hz trong chân không. Bước sóng tính theo nm bằng bao nhiêu?
\shortans{250}
\loigiai{
Xác định đúng dữ kiện và đơn vị, áp dụng hệ thức của dạng bài rồi tính được kết quả 250.
}
\end{ex}

% ===== Câu 88 =====
% ID: L11C2B11-23-TN
% Mức: TH
\dangbt{Nhận biết bản chất và tính chất của sóng điện từ}
\begin{ex}
% ID: L11C2B11-23-TN
% Mức: TH
Sóng điện từ có tốc độ $2\times10^8$ m/s trong môi trường. Trong 5 micro giây truyền được bao nhiêu m? Chọn giá trị đúng.
\choice
{2000}
{500}
{1001}
{\True 1000}
\loigiai{
Đáp án D. 1000. Đối chiếu định nghĩa, hệ thức hoặc dữ kiện của bài toán cho kết luận trên.
}
\end{ex}

% ===== Câu 89 =====
% ID: L11C2B11-24-DS
% Mức: VD
\dangbt{Tính bước sóng, tần số và tốc độ của sóng điện từ}
\begin{ex}
% ID: L11C2B11-24-DS
% Mức: VD
Xét các nhận định tổng hợp thuộc dạng “Tính bước sóng, tần số và tốc độ của sóng điện từ”.
\choiceTF
{\True $3\,\text{m}$.}
{\True $5\times10^{14}$ Hz.}
{\True $4\,\text{m}$.}
{$0,3\,\text{m}$.}
\loigiai{
\begin{itemchoice}
\item (Đúng) $3\,\text{m}$. (Vì): $3\,\text{m}$. b) Đ: $5\times10^{14}$ Hz. c) Đ: $4\,\text{m}$. d) S: $0,3\,\text{m}$
\item (Đúng) $5\times10^{14}$ Hz.
\item (Đúng) $4\,\text{m}$.
\item (Sai) $0,3\,\text{m}$.
\end{itemchoice}
}
\end{ex}

% ===== Câu 90 =====
% ID: L11C2B11-24-TL
% Mức: VD
\dangbt{Phân loại thang sóng điện từ và ứng dụng}
\begin{ex}
% ID: L11C2B11-24-TL
% Mức: VD
Giải thích vì sao cần hạn chế tiếp xúc tia tử ngoại và tia $X$ dù chúng có nhiều ứng dụng.
\choiceTF

\loigiai{
Đây là các bức xạ tần số cao, có thể gây tổn thương mô và vật liệu di truyền; cần kiểm soát liều và che chắn.
}
\end{ex}

% ===== Câu 91 =====
% ID: L11C2B11-24-TLN
% Mức: VD
\begin{ex}
% ID: L11C2B11-24-TLN
% Mức: VD
Từ kết quả của tình huống sau, nếu đại lượng cần tìm được nhân với hệ số 1.5 thì giá trị mới bằng bao nhiêu? Sóng vô tuyến có bước sóng $100\,\text{m}$. Tần số tính theo MHz bằng bao nhiêu?
\shortans{04/05/2026}
\loigiai{
Xác định đúng dữ kiện và đơn vị, áp dụng hệ thức của dạng bài rồi tính được kết quả 4.5.
}
\end{ex}

% ===== Câu 92 =====
% ID: L11C2B11-24-TN
% Mức: VD
\dangbt{Nhận biết bản chất và tính chất của sóng điện từ}
\begin{ex}
% ID: L11C2B11-24-TN
% Mức: VD
Kết luận nào phù hợp nhất khi giải bài toán sau: Nêu ba điểm khác nhau cơ bản giữa sóng điện từ và sóng cơ.
\choice
{\True Sóng điện từ là điện từ trường biến thiên, truyền được trong chân không và luôn là sóng ngang}
{Không cần sử dụng định nghĩa hay dữ kiện đã cho.}
{Đại lượng cần tìm luôn bằng 0 trong mọi trường hợp.}
{Kết quả không phụ thuộc môi trường, tần số hoặc điều kiện biên.}
\loigiai{
Đáp án A. Sóng điện từ là điện từ trường biến thiên, truyền được trong chân không và luôn là sóng ngang. Đối chiếu định nghĩa, hệ thức hoặc dữ kiện của bài toán cho kết luận trên.
}
\end{ex}

% ===== Câu 93 =====
% ID: L11C2B11-25-DS
% Mức: VD
\dangbt{Tính bước sóng, tần số và tốc độ của sóng điện từ}
\begin{ex}
% ID: L11C2B11-25-DS
% Mức: VD
Đánh giá cách xử lí các dữ kiện định lượng của dạng “Tính bước sóng, tần số và tốc độ của sóng điện từ”.
\choiceTF
{\True Kết quả của bài toán thứ nhất là 5.}
{\True Kết quả của bài toán thứ hai là 7.5.}
{\True Kết quả của bài toán thứ ba là 3.}
{Có thể bỏ qua hoàn toàn đơn vị và điều kiện vật lí khi tính toán.}
\loigiai{
\begin{itemchoice}
\item (Đúng) Kết quả của bài toán thứ nhất là 5. (Vì): Kết quả của bài toán thứ nhất là 5. b) Đ: Kết quả của bài toán thứ hai là 7.5. c) Đ: Kết quả của bài toán thứ ba là 3. d) S: Có thể bỏ qua hoàn toàn đơn vị và điều kiện vật lí khi tính toán
\item (Đúng) Kết quả của bài toán thứ hai là 7.5.
\item (Đúng) Kết quả của bài toán thứ ba là 3.
\item (Sai) Có thể bỏ qua hoàn toàn đơn vị và điều kiện vật lí khi tính toán.
\end{itemchoice}
}
\end{ex}

% ===== Câu 94 =====
% ID: L11C2B11-25-TL
% Mức: VD
\dangbt{Phân loại thang sóng điện từ và ứng dụng}
\begin{ex}
% ID: L11C2B11-25-TL
% Mức: VD
Phân tích các bước lập luận và chỉ ra điều kiện áp dụng trong bài toán: Lập thứ tự từ bước sóng dài đến ngắn cho: hồng ngoại, tử ngoại, sóng vô tuyến, ánh sáng nhìn thấy, tia X.
\choiceTF

\loigiai{
Sóng vô tuyến, hồng ngoại, ánh sáng nhìn thấy, tử ngoại, tia X. Cần nêu rõ điều kiện áp dụng, đổi đơn vị thống nhất và kiểm tra ý nghĩa vật lí của kết quả.
}
\end{ex}

% ===== Câu 95 =====
% ID: L11C2B11-25-TLN
% Mức: TH
\dangbt{Tính bước sóng, tần số và tốc độ của sóng điện từ}
\begin{ex}
% ID: L11C2B11-25-TLN
% Mức: TH
Sóng điện từ 60 MHz truyền trong chân không. Bước sóng tính theo m bằng bao nhiêu?
\shortans{5}
\loigiai{
Xác định đúng dữ kiện và đơn vị, áp dụng hệ thức của dạng bài rồi tính được kết quả 5.
}
\end{ex}

% ===== Câu 96 =====
% ID: L11C2B11-25-TN
% Mức: VD
\dangbt{Nhận biết bản chất và tính chất của sóng điện từ}
\begin{ex}
% ID: L11C2B11-25-TN
% Mức: VD
Kết luận nào phù hợp nhất khi giải bài toán sau: Mô tả quan hệ phương và pha của $\vec E$, $\vec B$ và phương truyền sóng điện từ.
\choice
{Kết quả không phụ thuộc môi trường, tần số hoặc điều kiện biên.}
{\True $\vec E$ vuông góc $\vec B$}
{Không cần sử dụng định nghĩa hay dữ kiện đã cho.}
{Đại lượng cần tìm luôn bằng 0 trong mọi trường hợp.}
\loigiai{
Đáp án B. $\vec E$ vuông góc $\vec B$. Đối chiếu định nghĩa, hệ thức hoặc dữ kiện của bài toán cho kết luận trên.
}
\end{ex}

% ===== Câu 97 =====
% ID: L11C2B11-26-DS
% Mức: TH
\dangbt{Vận dụng tính chất sóng điện từ trong truyền thông}
\begin{ex}
% ID: L11C2B11-26-DS
% Mức: TH
Trong chuyên đề “Vận dụng tính chất sóng điện từ trong truyền thông”, Xét sóng điện từ truyền trong chân không.
\choiceTF
{\True Sóng điện từ là sóng ngang.}
{\True Điện trường và từ trường biến thiên cùng pha.}
{\True Sóng mang năng lượng.}
{Sóng cần môi trường vật chất đàn hồi.}
\loigiai{
\begin{itemchoice}
\item (Đúng) Sóng điện từ là sóng ngang. (Vì): Sóng điện từ là sóng ngang. b) Đ: Điện trường và từ trường biến thiên cùng pha. c) Đ: Sóng mang năng lượng. d) S: Sóng cần môi trường vật chất đàn hồi
\item (Đúng) Điện trường và từ trường biến thiên cùng pha.
\item (Đúng) Sóng mang năng lượng.
\item (Sai) Sóng cần môi trường vật chất đàn hồi.
\end{itemchoice}
}
\end{ex}

% ===== Câu 98 =====
% ID: L11C2B11-26-TL
% Mức: VD
\dangbt{Phân loại thang sóng điện từ và ứng dụng}
\begin{ex}
% ID: L11C2B11-26-TL
% Mức: VD
Một học sinh giải bài sau nhưng bỏ qua điều kiện vật lí. Hãy sửa cách làm và trình bày lời giải đúng: Chọn bức xạ phù hợp cho điều khiển từ xa, chụp xương và truyền thông vệ tinh; giải thích ngắn gọn.
\choiceTF

\loigiai{
Cách giải đúng: Điều khiển từ xa: hồng ngoại; chụp xương: tia $X$ do khả năng đâm xuyên khác nhau; vệ tinh: vi ba do truyền định hướng và qua khí quyển phù hợp. Sai lầm cần tránh là dùng hệ thức khi chưa xác định điều kiện biên, môi trường hoặc đơn vị.
}
\end{ex}

% ===== Câu 99 =====
% ID: L11C2B11-26-TLN
% Mức: TH
\dangbt{Tính bước sóng, tần số và tốc độ của sóng điện từ}
\begin{ex}
% ID: L11C2B11-26-TLN
% Mức: TH
Bức xạ có bước sóng 400 nm trong chân không. Tần số tính theo $10^{14}$ Hz bằng bao nhiêu?
\shortans{07/05/2026}
\loigiai{
Xác định đúng dữ kiện và đơn vị, áp dụng hệ thức của dạng bài rồi tính được kết quả 7.5.
}
\end{ex}

% ===== Câu 100 =====
% ID: L11C2B11-26-TN
% Mức: VD
\dangbt{Nhận biết bản chất và tính chất của sóng điện từ}
\begin{ex}
% ID: L11C2B11-26-TN
% Mức: VD
Kết luận nào phù hợp nhất khi giải bài toán sau: Giải thích vì sao ánh sáng được xem là sóng điện từ.
\choice
{Đại lượng cần tìm luôn bằng 0 trong mọi trường hợp.}
{Kết quả không phụ thuộc môi trường, tần số hoặc điều kiện biên.}
{\True Ánh sáng có các tính chất sóng điện từ và truyền trong chân không với tốc độ đúng bằng tốc độ Maxwell dự đoán cho sóng điện từ.}
{Không cần sử dụng định nghĩa hay dữ kiện đã cho.}
\loigiai{
Đáp án C. Ánh sáng có các tính chất sóng điện từ và truyền trong chân không với tốc độ đúng bằng tốc độ Maxwell dự đoán cho sóng điện từ.. Đối chiếu định nghĩa, hệ thức hoặc dữ kiện của bài toán cho kết luận trên.
}
\end{ex}

% ===== Câu 101 =====
% ID: L11C2B11-27-DS
% Mức: TH
\dangbt{Vận dụng tính chất sóng điện từ trong truyền thông}
\begin{ex}
% ID: L11C2B11-27-DS
% Mức: TH
Trong chuyên đề “Vận dụng tính chất sóng điện từ trong truyền thông”, Sóng điện từ đi từ chân không vào điện môi.
\choiceTF
{\True Tần số không đổi.}
{\True Tốc độ thường giảm.}
{\True Bước sóng giảm nếu tốc độ giảm.}
{Chu kì tăng gấp đôi trong mọi điện môi.}
\loigiai{
\begin{itemchoice}
\item (Đúng) Tần số không đổi. (Vì): Tần số không đổi. b) Đ: Tốc độ thường giảm. c) Đ: Bước sóng giảm nếu tốc độ giảm. d) S: Chu kì tăng gấp đôi trong mọi điện môi
\item (Đúng) Tốc độ thường giảm.
\item (Đúng) Bước sóng giảm nếu tốc độ giảm.
\item (Sai) Chu kì tăng gấp đôi trong mọi điện môi.
\end{itemchoice}
}
\end{ex}

% ===== Câu 102 =====
% ID: L11C2B11-27-TL
% Mức: VD
\dangbt{Phân loại thang sóng điện từ và ứng dụng}
\begin{ex}
% ID: L11C2B11-27-TL
% Mức: VD
Mở rộng tình huống sau thành một ví dụ thực tế và trình bày phương pháp giải: Giải thích vì sao cần hạn chế tiếp xúc tia tử ngoại và tia $X$ dù chúng có nhiều ứng dụng.
\choiceTF

\loigiai{
Đây là các bức xạ tần số cao, có thể gây tổn thương mô và vật liệu di truyền; cần kiểm soát liều và che chắn. Có thể kiểm chứng bằng cách thay kết quả vào dữ kiện ban đầu và đối chiếu với hiện tượng thực tế.
}
\end{ex}

% ===== Câu 103 =====
% ID: L11C2B11-27-TLN
% Mức: TH
\dangbt{Tính bước sóng, tần số và tốc độ của sóng điện từ}
\begin{ex}
% ID: L11C2B11-27-TLN
% Mức: TH
Sóng có tần số 80 MHz truyền với tốc độ $2,4\times10^8$ m/s. Bước sóng tính theo m bằng bao nhiêu?
\shortans{3}
\loigiai{
Xác định đúng dữ kiện và đơn vị, áp dụng hệ thức của dạng bài rồi tính được kết quả 3.
}
\end{ex}

% ===== Câu 104 =====
% ID: L11C2B11-27-TN
% Mức: VD
\dangbt{Nhận biết bản chất và tính chất của sóng điện từ}
\begin{ex}
% ID: L11C2B11-27-TN
% Mức: VD
Phương pháp phù hợp nhất cho dạng “Nhận biết bản chất và tính chất của sóng điện từ” là
\choice
{chỉ thay số mà không cần xét điều kiện áp dụng}
{bỏ qua đơn vị và chiều truyền sóng}
{luôn coi mọi điểm dao động cùng pha}
{\True xác định dữ kiện, chọn đúng hệ thức hoặc định nghĩa, đổi đơn vị rồi kiểm tra kết quả}
\loigiai{
Đáp án D. xác định dữ kiện, chọn đúng hệ thức hoặc định nghĩa, đổi đơn vị rồi kiểm tra kết quả. Đối chiếu định nghĩa, hệ thức hoặc dữ kiện của bài toán cho kết luận trên.
}
\end{ex}

% ===== Câu 105 =====
% ID: L11C2B11-28-DS
% Mức: TH
\dangbt{Vận dụng tính chất sóng điện từ trong truyền thông}
\begin{ex}
% ID: L11C2B11-28-DS
% Mức: TH
Trong chuyên đề “Vận dụng tính chất sóng điện từ trong truyền thông”, Các hiện tượng của sóng điện từ.
\choiceTF
{\True Có thể phản xạ.}
{\True Có thể khúc xạ.}
{\True Có thể giao thoa và nhiễu xạ.}
{Không bao giờ truyền qua chân không.}
\loigiai{
\begin{itemchoice}
\item (Đúng) Có thể phản xạ. (Vì): Có thể phản xạ. b) Đ: Có thể khúc xạ. c) Đ: Có thể giao thoa và nhiễu xạ. d) S: Không bao giờ truyền qua chân không
\item (Đúng) Có thể khúc xạ.
\item (Đúng) Có thể giao thoa và nhiễu xạ.
\item (Sai) Không bao giờ truyền qua chân không.
\end{itemchoice}
}
\end{ex}

% ===== Câu 106 =====
% ID: L11C2B11-28-TL
% Mức: TH
\dangbt{Tính bước sóng, tần số và tốc độ của sóng điện từ}
\begin{ex}
% ID: L11C2B11-28-TL
% Mức: TH
Một đài phát ở tần số 100 MHz. Tính bước sóng trong chân không và trong môi trường có tốc độ $2,4\times10^8$ m/s.
\choiceTF

\loigiai{
Trong chân không $\lambda=3$ m; trong môi trường $\lambda=2,4$ m; tần số không đổi.
}
\end{ex}

% ===== Câu 107 =====
% ID: L11C2B11-28-TLN
% Mức: VD
\begin{ex}
% ID: L11C2B11-28-TLN
% Mức: VD
Từ kết quả của tình huống sau, nếu đại lượng cần tìm được nhân với hệ số 2 thì giá trị mới bằng bao nhiêu? Sóng điện từ 60 MHz truyền trong chân không. Bước sóng tính theo m bằng bao nhiêu?
\shortans{10}
\loigiai{
Xác định đúng dữ kiện và đơn vị, áp dụng hệ thức của dạng bài rồi tính được kết quả 10.
}
\end{ex}

% ===== Câu 108 =====
% ID: L11C2B11-28-TN
% Mức: NB
\dangbt{Nhận biết bức xạ điện từ theo nguồn phát và công dụng}
\begin{ex}
% ID: L11C2B11-28-TN
% Mức: NB
Trong chuyên đề “Nhận biết bức xạ điện từ theo nguồn phát và công dụng”, Trong thang sóng điện từ, bức xạ có bước sóng dài nhất trong các lựa chọn là
\choice
{tia $X$}
{tia gamma}
{\True sóng vô tuyến}
{tia tử ngoại}
\loigiai{
Đáp án C. sóng vô tuyến. Đối chiếu định nghĩa, hệ thức hoặc dữ kiện của bài toán cho kết luận trên.
}
\end{ex}

% ===== Câu 109 =====
% ID: L11C2B11-29-DS
% Mức: VD
\dangbt{Vận dụng tính chất sóng điện từ trong truyền thông}
\begin{ex}
% ID: L11C2B11-29-DS
% Mức: VD
Trong chuyên đề “Vận dụng tính chất sóng điện từ trong truyền thông”, Xét các nhận định tổng hợp thuộc dạng “Vận dụng tính chất sóng điện từ trong truyền thông”.
\choiceTF
{\True điện từ trường biến thiên.}
{\True truyền được với tốc độ $3\times10^8$ m/s.}
{\True vuông góc nhau và vuông góc phương truyền.}
{các phân tử mang điện đi xa.}
\loigiai{
\begin{itemchoice}
\item (Đúng) điện từ trường biến thiên. (Vì): iện từ trường biến thiên. b) Đ: truyền được với tốc độ $3\times10^8$ m/s. c) Đ: vuông góc nhau và vuông góc phương truyền. d) S: các phân tử mang điện đi xa
\item (Đúng) truyền được với tốc độ $3\times10^8$ m/s.
\item (Đúng) vuông góc nhau và vuông góc phương truyền.
\item (Sai) các phân tử mang điện đi xa.
\end{itemchoice}
}
\end{ex}

% ===== Câu 110 =====
% ID: L11C2B11-29-TL
% Mức: TH
\dangbt{Tính bước sóng, tần số và tốc độ của sóng điện từ}
\begin{ex}
% ID: L11C2B11-29-TL
% Mức: TH
Ánh sáng đỏ có bước sóng 750 nm, ánh sáng tím 400 nm. Tính tần số mỗi ánh sáng và so sánh.
\choiceTF

\loigiai{
$f_đ=4,0\times10^{14}$ Hz; $f_t=7,5\times10^{14}$ Hz; ánh sáng tím có tần số lớn hơn.
}
\end{ex}

% ===== Câu 111 =====
% ID: L11C2B11-29-TLN
% Mức: VD
\begin{ex}
% ID: L11C2B11-29-TLN
% Mức: VD
Từ kết quả của tình huống sau, nếu đại lượng cần tìm được nhân với hệ số 0.5 thì giá trị mới bằng bao nhiêu? Bức xạ có bước sóng 400 nm trong chân không. Tần số tính theo $10^{14}$ Hz bằng bao nhiêu?
\shortans{3.75}
\loigiai{
Xác định đúng dữ kiện và đơn vị, áp dụng hệ thức của dạng bài rồi tính được kết quả 3.75.
}
\end{ex}

% ===== Câu 112 =====
% ID: L11C2B11-29-TN
% Mức: NB
\dangbt{Nhận biết bức xạ điện từ theo nguồn phát và công dụng}
\begin{ex}
% ID: L11C2B11-29-TN
% Mức: NB
Trong chuyên đề “Nhận biết bức xạ điện từ theo nguồn phát và công dụng”, Bức xạ thường dùng trong điều khiển từ xa là
\choice
{tia gamma}
{tia $X$}
{tia tử ngoại}
{\True tia hồng ngoại}
\loigiai{
Đáp án D. tia hồng ngoại. Đối chiếu định nghĩa, hệ thức hoặc dữ kiện của bài toán cho kết luận trên.
}
\end{ex}

% ===== Câu 113 =====
% ID: L11C2B11-30-DS
% Mức: VD
\dangbt{Vận dụng tính chất sóng điện từ trong truyền thông}
\begin{ex}
% ID: L11C2B11-30-DS
% Mức: VD
Trong chuyên đề “Vận dụng tính chất sóng điện từ trong truyền thông”, Đánh giá cách xử lí các dữ kiện định lượng của dạng “Vận dụng tính chất sóng điện từ trong truyền thông”.
\choiceTF
{\True Kết quả của bài toán thứ nhất là 600.}
{\True Kết quả của bài toán thứ hai là 3.}
{\True Kết quả của bài toán thứ ba là 1000.}
{Có thể bỏ qua hoàn toàn đơn vị và điều kiện vật lí khi tính toán.}
\loigiai{
\begin{itemchoice}
\item (Đúng) Kết quả của bài toán thứ nhất là 600. (Vì): Kết quả của bài toán thứ nhất là 600. b) Đ: Kết quả của bài toán thứ hai là 3. c) Đ: Kết quả của bài toán thứ ba là 1000. d) S: Có thể bỏ qua hoàn toàn đơn vị và điều kiện vật lí khi tính toán
\item (Đúng) Kết quả của bài toán thứ hai là 3.
\item (Đúng) Kết quả của bài toán thứ ba là 1000.
\item (Sai) Có thể bỏ qua hoàn toàn đơn vị và điều kiện vật lí khi tính toán.
\end{itemchoice}
}
\end{ex}

% ===== Câu 114 =====
% ID: L11C2B11-30-TL
% Mức: VD
\dangbt{Tính bước sóng, tần số và tốc độ của sóng điện từ}
\begin{ex}
% ID: L11C2B11-30-TL
% Mức: VD
Một xung radar đi tới mục tiêu rồi phản xạ về sau 200 micro giây. Tính khoảng cách tới mục tiêu.
\choiceTF

\loigiai{
Quãng đường hai chiều $ct$ nên $d=ct/2=3\times10^8\times200\times10^{-6}/2=30000$ m = $30\,\text{km}$.
}
\end{ex}

% ===== Câu 115 =====
% ID: L11C2B11-30-TLN
% Mức: VD
\begin{ex}
% ID: L11C2B11-30-TLN
% Mức: VD
Từ kết quả của tình huống sau, nếu đại lượng cần tìm được nhân với hệ số 1.5 thì giá trị mới bằng bao nhiêu? Sóng có tần số 80 MHz truyền với tốc độ $2,4\times10^8$ m/s. Bước sóng tính theo m bằng bao nhiêu?
\shortans{04/05/2026}
\loigiai{
Xác định đúng dữ kiện và đơn vị, áp dụng hệ thức của dạng bài rồi tính được kết quả 4.5.
}
\end{ex}

% ===== Câu 116 =====
% ID: L11C2B11-30-TN
% Mức: TH
\dangbt{Nhận biết bức xạ điện từ theo nguồn phát và công dụng}
\begin{ex}
% ID: L11C2B11-30-TN
% Mức: TH
Trong chuyên đề “Nhận biết bức xạ điện từ theo nguồn phát và công dụng”, Tia $X$ được ứng dụng phổ biến để
\choice
{\True chụp ảnh xương}
{truyền thanh AM}
{sưởi bằng đèn nhiệt}
{điều khiển tivi}
\loigiai{
Đáp án A. chụp ảnh xương. Đối chiếu định nghĩa, hệ thức hoặc dữ kiện của bài toán cho kết luận trên.
}
\end{ex}

% ===== Câu 117 =====
% ID: L11C2B11-31-TL
% Mức: VD
\dangbt{Tính bước sóng, tần số và tốc độ của sóng điện từ}
\begin{ex}
% ID: L11C2B11-31-TL
% Mức: VD
Phân tích các bước lập luận và chỉ ra điều kiện áp dụng trong bài toán: Một đài phát ở tần số 100 MHz. Tính bước sóng trong chân không và trong môi trường có tốc độ $2,4\times10^8$ m/s.
\choiceTF

\loigiai{
Trong chân không $\lambda=3$ m; trong môi trường $\lambda=2,4$ m; tần số không đổi. Cần nêu rõ điều kiện áp dụng, đổi đơn vị thống nhất và kiểm tra ý nghĩa vật lí của kết quả.
}
\end{ex}

% ===== Câu 118 =====
% ID: L11C2B11-31-TLN
% Mức: TH
\dangbt{Vận dụng tính chất sóng điện từ trong truyền thông}
\begin{ex}
% ID: L11C2B11-31-TLN
% Mức: TH
Trong chuyên đề “Vận dụng tính chất sóng điện từ trong truyền thông”, Sóng điện từ truyền trong chân không $0,002\,\text{s}$. Quãng đường tính theo km bằng bao nhiêu?
\shortans{600}
\loigiai{
Xác định đúng dữ kiện và đơn vị, áp dụng hệ thức của dạng bài rồi tính được kết quả 600.
}
\end{ex}

% ===== Câu 119 =====
% ID: L11C2B11-31-TN
% Mức: TH
\dangbt{Nhận biết bức xạ điện từ theo nguồn phát và công dụng}
\begin{ex}
% ID: L11C2B11-31-TN
% Mức: TH
Trong chuyên đề “Nhận biết bức xạ điện từ theo nguồn phát và công dụng”, Bức xạ có bước sóng 600 nm. Đổi sang micromet bằng bao nhiêu? Chọn giá trị đúng.
\choice
{01/06/2026}
{\True 0.6}
{01/02/2026}
{0.3}
\loigiai{
Đáp án B. 0.6. Đối chiếu định nghĩa, hệ thức hoặc dữ kiện của bài toán cho kết luận trên.
}
\end{ex}

% ===== Câu 120 =====
% ID: L11C2B11-32-TL
% Mức: VD
\dangbt{Tính bước sóng, tần số và tốc độ của sóng điện từ}
\begin{ex}
% ID: L11C2B11-32-TL
% Mức: VD
Một học sinh giải bài sau nhưng bỏ qua điều kiện vật lí. Hãy sửa cách làm và trình bày lời giải đúng: Ánh sáng đỏ có bước sóng 750 nm, ánh sáng tím 400 nm. Tính tần số mỗi ánh sáng và so sánh.
\choiceTF

\loigiai{
Cách giải đúng: $f_đ=4,0\times10^{14}$ Hz; $f_t=7,5\times10^{14}$ Hz; ánh sáng tím có tần số lớn hơn. Sai lầm cần tránh là dùng hệ thức khi chưa xác định điều kiện biên, môi trường hoặc đơn vị.
}
\end{ex}

% ===== Câu 121 =====
% ID: L11C2B11-32-TLN
% Mức: TH
\dangbt{Vận dụng tính chất sóng điện từ trong truyền thông}
\begin{ex}
% ID: L11C2B11-32-TLN
% Mức: TH
Trong chuyên đề “Vận dụng tính chất sóng điện từ trong truyền thông”, Ánh sáng đi $900\,\text{km}$ trong chân không. Thời gian tính theo ms bằng bao nhiêu?
\shortans{3}
\loigiai{
Xác định đúng dữ kiện và đơn vị, áp dụng hệ thức của dạng bài rồi tính được kết quả 3.
}
\end{ex}

% ===== Câu 122 =====
% ID: L11C2B11-32-TN
% Mức: TH
\dangbt{Nhận biết bức xạ điện từ theo nguồn phát và công dụng}
\begin{ex}
% ID: L11C2B11-32-TN
% Mức: TH
Trong chuyên đề “Nhận biết bức xạ điện từ theo nguồn phát và công dụng”, Tia có tần số $6\times10^{14}$ Hz trong chân không. Bước sóng tính theo nm bằng bao nhiêu? Chọn giá trị đúng.
\choice
{250}
{501}
{\True 500}
{1000}
\loigiai{
Đáp án C. 500. Đối chiếu định nghĩa, hệ thức hoặc dữ kiện của bài toán cho kết luận trên.
}
\end{ex}

% ===== Câu 123 =====
% ID: L11C2B11-33-TL
% Mức: VD
\dangbt{Tính bước sóng, tần số và tốc độ của sóng điện từ}
\begin{ex}
% ID: L11C2B11-33-TL
% Mức: VD
Mở rộng tình huống sau thành một ví dụ thực tế và trình bày phương pháp giải: Một xung radar đi tới mục tiêu rồi phản xạ về sau 200 micro giây. Tính khoảng cách tới mục tiêu.
\choiceTF

\loigiai{
Quãng đường hai chiều $ct$ nên $d=ct/2=3\times10^8\times200\times10^{-6}/2=30000$ m = $30\,\text{km}$. Có thể kiểm chứng bằng cách thay kết quả vào dữ kiện ban đầu và đối chiếu với hiện tượng thực tế.
}
\end{ex}

% ===== Câu 124 =====
% ID: L11C2B11-33-TLN
% Mức: TH
\dangbt{Vận dụng tính chất sóng điện từ trong truyền thông}
\begin{ex}
% ID: L11C2B11-33-TLN
% Mức: TH
Trong chuyên đề “Vận dụng tính chất sóng điện từ trong truyền thông”, Sóng điện từ có tốc độ $2\times10^8$ m/s trong môi trường. Trong 5 micro giây truyền được bao nhiêu m?
\shortans{1000}
\loigiai{
Xác định đúng dữ kiện và đơn vị, áp dụng hệ thức của dạng bài rồi tính được kết quả 1000.
}
\end{ex}

% ===== Câu 125 =====
% ID: L11C2B11-33-TN
% Mức: TH
\dangbt{Nhận biết bức xạ điện từ theo nguồn phát và công dụng}
\begin{ex}
% ID: L11C2B11-33-TN
% Mức: TH
Trong chuyên đề “Nhận biết bức xạ điện từ theo nguồn phát và công dụng”, Sóng vô tuyến có bước sóng $100\,\text{m}$. Tần số tính theo MHz bằng bao nhiêu? Chọn giá trị đúng.
\choice
{6}
{1.5}
{4}
{\True 2.1}
\loigiai{
Đáp án D. 3. Đối chiếu định nghĩa, hệ thức hoặc dữ kiện của bài toán cho kết luận trên.
}
\end{ex}

% ===== Câu 126 =====
% ID: L11C2B11-34-TL
% Mức: TH
\dangbt{Vận dụng tính chất sóng điện từ trong truyền thông}
\begin{ex}
% ID: L11C2B11-34-TL
% Mức: TH
Trong chuyên đề “Vận dụng tính chất sóng điện từ trong truyền thông”, Nêu ba điểm khác nhau cơ bản giữa sóng điện từ và sóng cơ.
\choiceTF

\loigiai{
Sóng điện từ là điện từ trường biến thiên, truyền được trong chân không và luôn là sóng ngang; sóng cơ là dao động vật chất, cần môi trường và có thể ngang hoặc dọc.
}
\end{ex}

% ===== Câu 127 =====
% ID: L11C2B11-34-TLN
% Mức: VD
\begin{ex}
% ID: L11C2B11-34-TLN
% Mức: VD
Trong chuyên đề “Vận dụng tính chất sóng điện từ trong truyền thông”, Từ kết quả của tình huống sau, nếu đại lượng cần tìm được nhân với hệ số 2 thì giá trị mới bằng bao nhiêu? Sóng điện từ truyền trong chân không $0,002\,\text{s}$. Quãng đường tính theo km bằng bao nhiêu?
\shortans{1200}
\loigiai{
Xác định đúng dữ kiện và đơn vị, áp dụng hệ thức của dạng bài rồi tính được kết quả 1200.
}
\end{ex}

% ===== Câu 128 =====
% ID: L11C2B11-34-TN
% Mức: VD
\dangbt{Nhận biết bức xạ điện từ theo nguồn phát và công dụng}
\begin{ex}
% ID: L11C2B11-34-TN
% Mức: VD
Trong chuyên đề “Nhận biết bức xạ điện từ theo nguồn phát và công dụng”, Kết luận nào phù hợp nhất khi giải bài toán sau: Lập thứ tự từ bước sóng dài đến ngắn cho: hồng ngoại, tử ngoại, sóng vô tuyến, ánh sáng nhìn thấy, tia X.
\choice
{\True Sóng vô tuyến, hồng ngoại, ánh sáng nhìn thấy, tử ngoại, tia X.}
{Không cần sử dụng định nghĩa hay dữ kiện đã cho.}
{Đại lượng cần tìm luôn bằng 0 trong mọi trường hợp.}
{Kết quả không phụ thuộc môi trường, tần số hoặc điều kiện biên.}
\loigiai{
Đáp án A. Sóng vô tuyến, hồng ngoại, ánh sáng nhìn thấy, tử ngoại, tia X.. Đối chiếu định nghĩa, hệ thức hoặc dữ kiện của bài toán cho kết luận trên.
}
\end{ex}

% ===== Câu 129 =====
% ID: L11C2B11-35-TL
% Mức: TH
\dangbt{Vận dụng tính chất sóng điện từ trong truyền thông}
\begin{bt}
% ID: L11C2B11-35-TL
% Mức: TH
Trong chuyên đề “Vận dụng tính chất sóng điện từ trong truyền thông”, Mô tả quan hệ phương và pha của $\vec E$, $\vec B$ và phương truyền sóng điện từ.
\loigiai{
$\vec E$ vuông góc $\vec B$; cả hai vuông góc phương truyền và biến thiên cùng pha tại một điểm.
}
\end{bt}

% ===== Câu 130 =====
% ID: L11C2B11-35-TLN
% Mức: VD
\begin{ex}
% ID: L11C2B11-35-TLN
% Mức: VD
Trong chuyên đề “Vận dụng tính chất sóng điện từ trong truyền thông”, Từ kết quả của tình huống sau, nếu đại lượng cần tìm được nhân với hệ số 0.5 thì giá trị mới bằng bao nhiêu? Ánh sáng đi $900\,\text{km}$ trong chân không. Thời gian tính theo ms bằng bao nhiêu?
\shortans{01/05/2026}
\loigiai{
Xác định đúng dữ kiện và đơn vị, áp dụng hệ thức của dạng bài rồi tính được kết quả 1.5.
}
\end{ex}

% ===== Câu 131 =====
% ID: L11C2B11-35-TN
% Mức: VD
\dangbt{Nhận biết bức xạ điện từ theo nguồn phát và công dụng}
\begin{ex}
% ID: L11C2B11-35-TN
% Mức: VD
Trong chuyên đề “Nhận biết bức xạ điện từ theo nguồn phát và công dụng”, Kết luận nào phù hợp nhất khi giải bài toán sau: Chọn bức xạ phù hợp cho điều khiển từ xa, chụp xương và truyền thông vệ tinh; giải thích ngắn gọn.
\choice
{Kết quả không phụ thuộc môi trường, tần số hoặc điều kiện biên.}
{\True Điều khiển từ xa: hồng ngoại}
{Không cần sử dụng định nghĩa hay dữ kiện đã cho.}
{Đại lượng cần tìm luôn bằng 0 trong mọi trường hợp.}
\loigiai{
Đáp án B. Điều khiển từ xa: hồng ngoại. Đối chiếu định nghĩa, hệ thức hoặc dữ kiện của bài toán cho kết luận trên.
}
\end{ex}

% ===== Câu 132 =====
% ID: L11C2B11-36-TL
% Mức: VD
\dangbt{Vận dụng tính chất sóng điện từ trong truyền thông}
\begin{ex}
% ID: L11C2B11-36-TL
% Mức: VD
Trong chuyên đề “Vận dụng tính chất sóng điện từ trong truyền thông”, Giải thích vì sao ánh sáng được xem là sóng điện từ.
\choiceTF

\loigiai{
Ánh sáng có các tính chất sóng điện từ và truyền trong chân không với tốc độ đúng bằng tốc độ Maxwell dự đoán cho sóng điện từ.
}
\end{ex}

% ===== Câu 133 =====
% ID: L11C2B11-36-TLN
% Mức: VD
\begin{ex}
% ID: L11C2B11-36-TLN
% Mức: VD
Trong chuyên đề “Vận dụng tính chất sóng điện từ trong truyền thông”, Từ kết quả của tình huống sau, nếu đại lượng cần tìm được nhân với hệ số 1.5 thì giá trị mới bằng bao nhiêu? Sóng điện từ có tốc độ $2\times10^8$ m/s trong môi trường. Trong 5 micro giây truyền được bao nhiêu m?
\shortans{1500}
\loigiai{
Xác định đúng dữ kiện và đơn vị, áp dụng hệ thức của dạng bài rồi tính được kết quả 1500.
}
\end{ex}

% ===== Câu 134 =====
% ID: L11C2B11-36-TN
% Mức: VD
\dangbt{Nhận biết bức xạ điện từ theo nguồn phát và công dụng}
\begin{ex}
% ID: L11C2B11-36-TN
% Mức: VD
Trong chuyên đề “Nhận biết bức xạ điện từ theo nguồn phát và công dụng”, Kết luận nào phù hợp nhất khi giải bài toán sau: Giải thích vì sao cần hạn chế tiếp xúc tia tử ngoại và tia $X$ dù chúng có nhiều ứng dụng.
\choice
{Đại lượng cần tìm luôn bằng 0 trong mọi trường hợp.}
{Kết quả không phụ thuộc môi trường, tần số hoặc điều kiện biên.}
{\True Đây là các bức xạ tần số cao, có thể gây tổn thương mô và vật liệu di truyền}
{Không cần sử dụng định nghĩa hay dữ kiện đã cho.}
\loigiai{
Đáp án C. Đây là các bức xạ tần số cao, có thể gây tổn thương mô và vật liệu di truyền. Đối chiếu định nghĩa, hệ thức hoặc dữ kiện của bài toán cho kết luận trên.
}
\end{ex}

% ===== Câu 135 =====
% ID: L11C2B11-37-TL
% Mức: VD
\dangbt{Vận dụng tính chất sóng điện từ trong truyền thông}
\begin{ex}
% ID: L11C2B11-37-TL
% Mức: VD
Trong chuyên đề “Vận dụng tính chất sóng điện từ trong truyền thông”, Phân tích các bước lập luận và chỉ ra điều kiện áp dụng trong bài toán: Nêu ba điểm khác nhau cơ bản giữa sóng điện từ và sóng cơ.
\choiceTF

\loigiai{
Sóng điện từ là điện từ trường biến thiên, truyền được trong chân không và luôn là sóng ngang; sóng cơ là dao động vật chất, cần môi trường và có thể ngang hoặc dọc. Cần nêu rõ điều kiện áp dụng, đổi đơn vị thống nhất và kiểm tra ý nghĩa vật lí của kết quả.
}
\end{ex}

% ===== Câu 136 =====
% ID: L11C2B11-37-TN
% Mức: VD
\dangbt{Nhận biết bức xạ điện từ theo nguồn phát và công dụng}
\begin{ex}
% ID: L11C2B11-37-TN
% Mức: VD
Trong chuyên đề “Nhận biết bức xạ điện từ theo nguồn phát và công dụng”, Phương pháp phù hợp nhất cho dạng “Nhận biết bức xạ điện từ theo nguồn phát và công dụng” là
\choice
{chỉ thay số mà không cần xét điều kiện áp dụng}
{bỏ qua đơn vị và chiều truyền sóng}
{luôn coi mọi điểm dao động cùng pha}
{\True xác định dữ kiện, chọn đúng hệ thức hoặc định nghĩa, đổi đơn vị rồi kiểm tra kết quả}
\loigiai{
Đáp án D. xác định dữ kiện, chọn đúng hệ thức hoặc định nghĩa, đổi đơn vị rồi kiểm tra kết quả. Đối chiếu định nghĩa, hệ thức hoặc dữ kiện của bài toán cho kết luận trên.
}
\end{ex}

% ===== Câu 137 =====
% ID: L11C2B11-38-TL
% Mức: VD
\dangbt{Vận dụng tính chất sóng điện từ trong truyền thông}
\begin{ex}
% ID: L11C2B11-38-TL
% Mức: VD
Trong chuyên đề “Vận dụng tính chất sóng điện từ trong truyền thông”, Một học sinh giải bài sau nhưng bỏ qua điều kiện vật lí. Hãy sửa cách làm và trình bày lời giải đúng: Mô tả quan hệ phương và pha của $\vec E$, $\vec B$ và phương truyền sóng điện từ.
\choiceTF

\loigiai{
Cách giải đúng: $\vec E$ vuông góc $\vec B$; cả hai vuông góc phương truyền và biến thiên cùng pha tại một điểm. Sai lầm cần tránh là dùng hệ thức khi chưa xác định điều kiện biên, môi trường hoặc đơn vị.
}
\end{ex}

% ===== Câu 138 =====
% ID: L11C2B11-38-TN
% Mức: NB
\dangbt{Phân loại thang sóng điện từ và ứng dụng}
\begin{ex}
% ID: L11C2B11-38-TN
% Mức: NB
Trong thang sóng điện từ, bức xạ có bước sóng dài nhất trong các lựa chọn là
\choice
{tia tử ngoại}
{tia $X$}
{tia gamma}
{\True sóng vô tuyến}
\loigiai{
Đáp án D. sóng vô tuyến. Đối chiếu định nghĩa, hệ thức hoặc dữ kiện của bài toán cho kết luận trên.
}
\end{ex}

% ===== Câu 139 =====
% ID: L11C2B11-39-TL
% Mức: VD
\dangbt{Vận dụng tính chất sóng điện từ trong truyền thông}
\begin{ex}
% ID: L11C2B11-39-TL
% Mức: VD
Trong chuyên đề “Vận dụng tính chất sóng điện từ trong truyền thông”, Mở rộng tình huống sau thành một ví dụ thực tế và trình bày phương pháp giải: Giải thích vì sao ánh sáng được xem là sóng điện từ.
\choiceTF

\loigiai{
Ánh sáng có các tính chất sóng điện từ và truyền trong chân không với tốc độ đúng bằng tốc độ Maxwell dự đoán cho sóng điện từ. Có thể kiểm chứng bằng cách thay kết quả vào dữ kiện ban đầu và đối chiếu với hiện tượng thực tế.
}
\end{ex}

% ===== Câu 140 =====
% ID: L11C2B11-39-TN
% Mức: NB
\dangbt{Phân loại thang sóng điện từ và ứng dụng}
\begin{ex}
% ID: L11C2B11-39-TN
% Mức: NB
Bức xạ thường dùng trong điều khiển từ xa là
\choice
{\True tia hồng ngoại}
{tia gamma}
{tia $X$}
{tia tử ngoại}
\loigiai{
Đáp án A. tia hồng ngoại. Đối chiếu định nghĩa, hệ thức hoặc dữ kiện của bài toán cho kết luận trên.
}
\end{ex}

% ===== Câu 141 =====
% ID: L11C2B11-40-TN
% Mức: TH
\begin{ex}
% ID: L11C2B11-40-TN
% Mức: TH
Tia $X$ được ứng dụng phổ biến để
\choice
{điều khiển tivi}
{\True chụp ảnh xương}
{truyền thanh AM}
{sưởi bằng đèn nhiệt}
\loigiai{
Đáp án B. chụp ảnh xương. Đối chiếu định nghĩa, hệ thức hoặc dữ kiện của bài toán cho kết luận trên.
}
\end{ex}

% ===== Câu 142 =====
% ID: L11C2B11-41-TN
% Mức: TH
\begin{ex}
% ID: L11C2B11-41-TN
% Mức: TH
Bức xạ có bước sóng 600 nm. Đổi sang micromet bằng bao nhiêu? Chọn giá trị đúng.
\choice
{0.3}
{01/06/2026}
{\True 0.6}
{01/02/2026}
\loigiai{
Đáp án C. 0.6. Đối chiếu định nghĩa, hệ thức hoặc dữ kiện của bài toán cho kết luận trên.
}
\end{ex}

% ===== Câu 143 =====
% ID: L11C2B11-42-TN
% Mức: TH
\begin{ex}
% ID: L11C2B11-42-TN
% Mức: TH
Tia có tần số $6\times10^{14}$ Hz trong chân không. Bước sóng tính theo nm bằng bao nhiêu? Chọn giá trị đúng.
\choice
{1000}
{250}
{501}
{\True 500}
\loigiai{
Đáp án D. 500. Đối chiếu định nghĩa, hệ thức hoặc dữ kiện của bài toán cho kết luận trên.
}
\end{ex}

% ===== Câu 144 =====
% ID: L11C2B11-43-TN
% Mức: TH
\begin{ex}
% ID: L11C2B11-43-TN
% Mức: TH
Sóng vô tuyến có bước sóng $100\,\text{m}$. Tần số tính theo MHz bằng bao nhiêu? Chọn giá trị đúng.
\choice
{\True 3}
{6}
{01/05/2026}
{4}
\loigiai{
Đáp án A. 3. Đối chiếu định nghĩa, hệ thức hoặc dữ kiện của bài toán cho kết luận trên.
}
\end{ex}

% ===== Câu 145 =====
% ID: L11C2B11-44-TN
% Mức: VD
\begin{ex}
% ID: L11C2B11-44-TN
% Mức: VD
Kết luận nào phù hợp nhất khi giải bài toán sau: Lập thứ tự từ bước sóng dài đến ngắn cho: hồng ngoại, tử ngoại, sóng vô tuyến, ánh sáng nhìn thấy, tia X.
\choice
{Kết quả không phụ thuộc môi trường, tần số hoặc điều kiện biên.}
{\True Sóng vô tuyến, hồng ngoại, ánh sáng nhìn thấy, tử ngoại, tia X.}
{Không cần sử dụng định nghĩa hay dữ kiện đã cho.}
{Đại lượng cần tìm luôn bằng 0 trong mọi trường hợp.}
\loigiai{
Đáp án B. Sóng vô tuyến, hồng ngoại, ánh sáng nhìn thấy, tử ngoại, tia X.. Đối chiếu định nghĩa, hệ thức hoặc dữ kiện của bài toán cho kết luận trên.
}
\end{ex}

% ===== Câu 146 =====
% ID: L11C2B11-45-TN
% Mức: VD
\begin{ex}
% ID: L11C2B11-45-TN
% Mức: VD
Kết luận nào phù hợp nhất khi giải bài toán sau: Chọn bức xạ phù hợp cho điều khiển từ xa, chụp xương và truyền thông vệ tinh; giải thích ngắn gọn.
\choice
{Đại lượng cần tìm luôn bằng 0 trong mọi trường hợp.}
{Kết quả không phụ thuộc môi trường, tần số hoặc điều kiện biên.}
{\True Điều khiển từ xa: hồng ngoại}
{Không cần sử dụng định nghĩa hay dữ kiện đã cho.}
\loigiai{
Đáp án C. Điều khiển từ xa: hồng ngoại. Đối chiếu định nghĩa, hệ thức hoặc dữ kiện của bài toán cho kết luận trên.
}
\end{ex}

% ===== Câu 147 =====
% ID: L11C2B11-46-TN
% Mức: VD
\begin{ex}
% ID: L11C2B11-46-TN
% Mức: VD
Kết luận nào phù hợp nhất khi giải bài toán sau: Giải thích vì sao cần hạn chế tiếp xúc tia tử ngoại và tia $X$ dù chúng có nhiều ứng dụng.
\choice
{Không cần sử dụng định nghĩa hay dữ kiện đã cho.}
{Đại lượng cần tìm luôn bằng 0 trong mọi trường hợp.}
{Kết quả không phụ thuộc môi trường, tần số hoặc điều kiện biên.}
{\True Đây là các bức xạ tần số cao, có thể gây tổn thương mô và vật liệu di truyền}
\loigiai{
Đáp án D. Đây là các bức xạ tần số cao, có thể gây tổn thương mô và vật liệu di truyền. Đối chiếu định nghĩa, hệ thức hoặc dữ kiện của bài toán cho kết luận trên.
}
\end{ex}

% ===== Câu 148 =====
% ID: L11C2B11-47-TN
% Mức: VD
\begin{ex}
% ID: L11C2B11-47-TN
% Mức: VD
Phương pháp phù hợp nhất cho dạng “Phân loại thang sóng điện từ và ứng dụng” là
\choice
{\True xác định dữ kiện, chọn đúng hệ thức hoặc định nghĩa, đổi đơn vị rồi kiểm tra kết quả}
{chỉ thay số mà không cần xét điều kiện áp dụng}
{bỏ qua đơn vị và chiều truyền sóng}
{luôn coi mọi điểm dao động cùng pha}
\loigiai{
Đáp án A. xác định dữ kiện, chọn đúng hệ thức hoặc định nghĩa, đổi đơn vị rồi kiểm tra kết quả. Đối chiếu định nghĩa, hệ thức hoặc dữ kiện của bài toán cho kết luận trên.
}
\end{ex}

% ===== Câu 149 =====
% ID: L11C2B11-48-TN
% Mức: NB
\dangbt{Tính bước sóng, tần số và tốc độ của sóng điện từ}
\begin{ex}
% ID: L11C2B11-48-TN
% Mức: NB
Sóng điện từ trong chân không có tần số 100 MHz. Bước sóng bằng
\choice
{\True $3\,\text{m}$}
{$0,3\,\text{m}$}
{$30\,\text{m}$}
{$300\,\text{m}$}
\loigiai{
Đáp án A. $3\,\text{m}$. Đối chiếu định nghĩa, hệ thức hoặc dữ kiện của bài toán cho kết luận trên.
}
\end{ex}

% ===== Câu 150 =====
% ID: L11C2B11-49-TN
% Mức: NB
\begin{ex}
% ID: L11C2B11-49-TN
% Mức: NB
Bức xạ có bước sóng 600 nm trong chân không có tần số
\choice
{$5\times10^5$ Hz}
{\True $5\times10^{14}$ Hz}
{$2\times10^{14}$ Hz}
{$1,8\times10^{17}$ Hz}
\loigiai{
Đáp án B. $5\times10^{14}$ Hz. Đối chiếu định nghĩa, hệ thức hoặc dữ kiện của bài toán cho kết luận trên.
}
\end{ex}

% ===== Câu 151 =====
% ID: L11C2B11-50-TN
% Mức: TH
\begin{ex}
% ID: L11C2B11-50-TN
% Mức: TH
Khi sóng điện từ đi vào môi trường có tốc độ $2\times10^8$ m/s, tần số 50 MHz, bước sóng là
\choice
{$2,5\,\text{m}$}
{$0,25\,\text{m}$}
{\True $4\,\text{m}$}
{$6\,\text{m}$}
\loigiai{
Đáp án C. $4\,\text{m}$. Đối chiếu định nghĩa, hệ thức hoặc dữ kiện của bài toán cho kết luận trên.
}
\end{ex}

% ===== Câu 152 =====
% ID: L11C2B11-51-TN
% Mức: TH
\begin{ex}
% ID: L11C2B11-51-TN
% Mức: TH
Sóng điện từ 60 MHz truyền trong chân không. Bước sóng tính theo m bằng bao nhiêu? Chọn giá trị đúng.
\choice
{10}
{02/05/2026}
{6}
{\True 5}
\loigiai{
Đáp án D. 5. Đối chiếu định nghĩa, hệ thức hoặc dữ kiện của bài toán cho kết luận trên.
}
\end{ex}

% ===== Câu 153 =====
% ID: L11C2B11-52-TN
% Mức: TH
\begin{ex}
% ID: L11C2B11-52-TN
% Mức: TH
Bức xạ có bước sóng 400 nm trong chân không. Tần số tính theo $10^{14}$ Hz bằng bao nhiêu? Chọn giá trị đúng.
\choice
{\True 07/05/2026}
{15}
{3.75}
{08/05/2026}
\loigiai{
Đáp án A. 7.5. Đối chiếu định nghĩa, hệ thức hoặc dữ kiện của bài toán cho kết luận trên.
}
\end{ex}

% ===== Câu 154 =====
% ID: L11C2B11-53-TN
% Mức: TH
\begin{ex}
% ID: L11C2B11-53-TN
% Mức: TH
Sóng có tần số 80 MHz truyền với tốc độ $2,4\times10^8$ m/s. Bước sóng tính theo m bằng bao nhiêu? Chọn giá trị đúng.
\choice
{4}
{\True 3}
{6}
{01/05/2026}
\loigiai{
Đáp án B. 3. Đối chiếu định nghĩa, hệ thức hoặc dữ kiện của bài toán cho kết luận trên.
}
\end{ex}

% ===== Câu 155 =====
% ID: L11C2B11-54-TN
% Mức: VD
\begin{ex}
% ID: L11C2B11-54-TN
% Mức: VD
Kết luận nào phù hợp nhất khi giải bài toán sau: Một đài phát ở tần số 100 MHz. Tính bước sóng trong chân không và trong môi trường có tốc độ $2,4\times10^8$ m/s.
\choice
{Đại lượng cần tìm luôn bằng 0 trong mọi trường hợp.}
{Kết quả không phụ thuộc môi trường, tần số hoặc điều kiện biên.}
{\True Trong chân không $\lambda=3$ m}
{Không cần sử dụng định nghĩa hay dữ kiện đã cho.}
\loigiai{
Đáp án C. Trong chân không $\lambda=3$ m. Đối chiếu định nghĩa, hệ thức hoặc dữ kiện của bài toán cho kết luận trên.
}
\end{ex}

% ===== Câu 156 =====
% ID: L11C2B11-55-TN
% Mức: VD
\begin{ex}
% ID: L11C2B11-55-TN
% Mức: VD
Kết luận nào phù hợp nhất khi giải bài toán sau: Ánh sáng đỏ có bước sóng 750 nm, ánh sáng tím 400 nm. Tính tần số mỗi ánh sáng và so sánh.
\choice
{Không cần sử dụng định nghĩa hay dữ kiện đã cho.}
{Đại lượng cần tìm luôn bằng 0 trong mọi trường hợp.}
{Kết quả không phụ thuộc môi trường, tần số hoặc điều kiện biên.}
{\True $f_đ=4,0\times10^{14}$ Hz}
\loigiai{
Đáp án D. $f_đ=4,0\times10^{14}$ Hz. Đối chiếu định nghĩa, hệ thức hoặc dữ kiện của bài toán cho kết luận trên.
}
\end{ex}

% ===== Câu 157 =====
% ID: L11C2B11-56-TN
% Mức: VD
\begin{ex}
% ID: L11C2B11-56-TN
% Mức: VD
Kết luận nào phù hợp nhất khi giải bài toán sau: Một xung radar đi tới mục tiêu rồi phản xạ về sau 200 micro giây. Tính khoảng cách tới mục tiêu.
\choice
{\True Quãng đường hai chiều $ct$ nên $d=ct/2=3\times10^8\times200\times10^{-6}/2=30000$ m = $30\,\text{km}$.}
{Không cần sử dụng định nghĩa hay dữ kiện đã cho.}
{Đại lượng cần tìm luôn bằng 0 trong mọi trường hợp.}
{Kết quả không phụ thuộc môi trường, tần số hoặc điều kiện biên.}
\loigiai{
Đáp án A. Quãng đường hai chiều $ct$ nên $d=ct/2=3\times10^8\times200\times10^{-6}/2=30000$ m = $30\,\text{km}$.. Đối chiếu định nghĩa, hệ thức hoặc dữ kiện của bài toán cho kết luận trên.
}
\end{ex}

% ===== Câu 158 =====
% ID: L11C2B11-57-TN
% Mức: VD
\begin{ex}
% ID: L11C2B11-57-TN
% Mức: VD
Phương pháp phù hợp nhất cho dạng “Tính bước sóng, tần số và tốc độ của sóng điện từ” là
\choice
{luôn coi mọi điểm dao động cùng pha}
{\True xác định dữ kiện, chọn đúng hệ thức hoặc định nghĩa, đổi đơn vị rồi kiểm tra kết quả}
{chỉ thay số mà không cần xét điều kiện áp dụng}
{bỏ qua đơn vị và chiều truyền sóng}
\loigiai{
Đáp án B. xác định dữ kiện, chọn đúng hệ thức hoặc định nghĩa, đổi đơn vị rồi kiểm tra kết quả. Đối chiếu định nghĩa, hệ thức hoặc dữ kiện của bài toán cho kết luận trên.
}
\end{ex}

% ===== Câu 159 =====
% ID: L11C2B11-58-TN
% Mức: NB
\dangbt{Vận dụng tính chất sóng điện từ trong truyền thông}
\begin{ex}
% ID: L11C2B11-58-TN
% Mức: NB
Trong chuyên đề “Vận dụng tính chất sóng điện từ trong truyền thông”, Sóng điện từ là sự lan truyền của
\choice
{dòng điện không đổi}
{\True điện từ trường biến thiên}
{các phân tử mang điện đi xa}
{áp suất cơ học}
\loigiai{
Đáp án B. điện từ trường biến thiên. Đối chiếu định nghĩa, hệ thức hoặc dữ kiện của bài toán cho kết luận trên.
}
\end{ex}

% ===== Câu 160 =====
% ID: L11C2B11-59-TN
% Mức: NB
\begin{ex}
% ID: L11C2B11-59-TN
% Mức: NB
Trong chuyên đề “Vận dụng tính chất sóng điện từ trong truyền thông”, Trong chân không, sóng điện từ
\choice
{chỉ là sóng dọc}
{có tốc độ phụ thuộc tần số rất mạnh}
{\True truyền được với tốc độ $3\times10^8$ m/s}
{không truyền được}
\loigiai{
Đáp án C. truyền được với tốc độ $3\times10^8$ m/s. Đối chiếu định nghĩa, hệ thức hoặc dữ kiện của bài toán cho kết luận trên.
}
\end{ex}

% ===== Câu 161 =====
% ID: L11C2B11-60-TN
% Mức: TH
\begin{ex}
% ID: L11C2B11-60-TN
% Mức: TH
Trong chuyên đề “Vận dụng tính chất sóng điện từ trong truyền thông”, Trong sóng điện từ, vectơ điện trường và vectơ từ trường
\choice
{song song nhau}
{cùng phương truyền}
{luôn bằng 0}
{\True vuông góc nhau và vuông góc phương truyền}
\loigiai{
Đáp án D. vuông góc nhau và vuông góc phương truyền. Đối chiếu định nghĩa, hệ thức hoặc dữ kiện của bài toán cho kết luận trên.
}
\end{ex}

% ===== Câu 162 =====
% ID: L11C2B11-61-TN
% Mức: TH
\begin{ex}
% ID: L11C2B11-61-TN
% Mức: TH
Trong chuyên đề “Vận dụng tính chất sóng điện từ trong truyền thông”, Sóng điện từ truyền trong chân không $0,002\,\text{s}$. Quãng đường tính theo km bằng bao nhiêu? Chọn giá trị đúng.
\choice
{\True 600}
{1200}
{300}
{601}
\loigiai{
Đáp án A. 600. Đối chiếu định nghĩa, hệ thức hoặc dữ kiện của bài toán cho kết luận trên.
}
\end{ex}

% ===== Câu 163 =====
% ID: L11C2B11-62-TN
% Mức: TH
\begin{ex}
% ID: L11C2B11-62-TN
% Mức: TH
Trong chuyên đề “Vận dụng tính chất sóng điện từ trong truyền thông”, Ánh sáng đi $900\,\text{km}$ trong chân không. Thời gian tính theo ms bằng bao nhiêu? Chọn giá trị đúng.
\choice
{4}
{\True 3}
{6}
{01/05/2026}
\loigiai{
Đáp án B. 3. Đối chiếu định nghĩa, hệ thức hoặc dữ kiện của bài toán cho kết luận trên.
}
\end{ex}

% ===== Câu 164 =====
% ID: L11C2B11-63-TN
% Mức: TH
\begin{ex}
% ID: L11C2B11-63-TN
% Mức: TH
Trong chuyên đề “Vận dụng tính chất sóng điện từ trong truyền thông”, Sóng điện từ có tốc độ $2\times10^8$ m/s trong môi trường. Trong 5 micro giây truyền được bao nhiêu m? Chọn giá trị đúng.
\choice
{500}
{1001}
{\True 1000}
{2000}
\loigiai{
Đáp án C. 1000. Đối chiếu định nghĩa, hệ thức hoặc dữ kiện của bài toán cho kết luận trên.
}
\end{ex}

% ===== Câu 165 =====
% ID: L11C2B11-64-TN
% Mức: VD
\begin{ex}
% ID: L11C2B11-64-TN
% Mức: VD
Trong chuyên đề “Vận dụng tính chất sóng điện từ trong truyền thông”, Kết luận nào phù hợp nhất khi giải bài toán sau: Nêu ba điểm khác nhau cơ bản giữa sóng điện từ và sóng cơ.
\choice
{Không cần sử dụng định nghĩa hay dữ kiện đã cho.}
{Đại lượng cần tìm luôn bằng 0 trong mọi trường hợp.}
{Kết quả không phụ thuộc môi trường, tần số hoặc điều kiện biên.}
{\True Sóng điện từ là điện từ trường biến thiên, truyền được trong chân không và luôn là sóng ngang}
\loigiai{
Đáp án D. Sóng điện từ là điện từ trường biến thiên, truyền được trong chân không và luôn là sóng ngang. Đối chiếu định nghĩa, hệ thức hoặc dữ kiện của bài toán cho kết luận trên.
}
\end{ex}

% ===== Câu 166 =====
% ID: L11C2B11-65-TN
% Mức: VD
\begin{ex}
% ID: L11C2B11-65-TN
% Mức: VD
Trong chuyên đề “Vận dụng tính chất sóng điện từ trong truyền thông”, Kết luận nào phù hợp nhất khi giải bài toán sau: Mô tả quan hệ phương và pha của $\vec E$, $\vec B$ và phương truyền sóng điện từ.
\choice
{\True $\vec E$ vuông góc $\vec B$}
{Không cần sử dụng định nghĩa hay dữ kiện đã cho.}
{Đại lượng cần tìm luôn bằng 0 trong mọi trường hợp.}
{Kết quả không phụ thuộc môi trường, tần số hoặc điều kiện biên.}
\loigiai{
Đáp án A. $\vec E$ vuông góc $\vec B$. Đối chiếu định nghĩa, hệ thức hoặc dữ kiện của bài toán cho kết luận trên.
}
\end{ex}

% ===== Câu 167 =====
% ID: L11C2B11-66-TN
% Mức: VD
\begin{ex}
% ID: L11C2B11-66-TN
% Mức: VD
Trong chuyên đề “Vận dụng tính chất sóng điện từ trong truyền thông”, Kết luận nào phù hợp nhất khi giải bài toán sau: Giải thích vì sao ánh sáng được xem là sóng điện từ.
\choice
{Kết quả không phụ thuộc môi trường, tần số hoặc điều kiện biên.}
{\True Ánh sáng có các tính chất sóng điện từ và truyền trong chân không với tốc độ đúng bằng tốc độ Maxwell dự đoán cho sóng điện từ.}
{Không cần sử dụng định nghĩa hay dữ kiện đã cho.}
{Đại lượng cần tìm luôn bằng 0 trong mọi trường hợp.}
\loigiai{
Đáp án B. Ánh sáng có các tính chất sóng điện từ và truyền trong chân không với tốc độ đúng bằng tốc độ Maxwell dự đoán cho sóng điện từ.. Đối chiếu định nghĩa, hệ thức hoặc dữ kiện của bài toán cho kết luận trên.
}
\end{ex}

% ===== Câu 168 =====
% ID: L11C2B11-67-TN
% Mức: VD
\begin{ex}
% ID: L11C2B11-67-TN
% Mức: VD
Trong chuyên đề “Vận dụng tính chất sóng điện từ trong truyền thông”, Phương pháp phù hợp nhất cho dạng “Vận dụng tính chất sóng điện từ trong truyền thông” là
\choice
{bỏ qua đơn vị và chiều truyền sóng}
{luôn coi mọi điểm dao động cùng pha}
{\True xác định dữ kiện, chọn đúng hệ thức hoặc định nghĩa, đổi đơn vị rồi kiểm tra kết quả}
{chỉ thay số mà không cần xét điều kiện áp dụng}
\loigiai{
Đáp án C. xác định dữ kiện, chọn đúng hệ thức hoặc định nghĩa, đổi đơn vị rồi kiểm tra kết quả. Đối chiếu định nghĩa, hệ thức hoặc dữ kiện của bài toán cho kết luận trên.
}
\end{ex}
